% 1.2 poskyris „Duomenų srautų diagramos“ (atsakingas Marijus Plančiūnas).
% Reikalingi paketai: tikz (bibliotekos arrows.meta, positioning, calc),
% float, pdflscape, longtable, booktabs, array.
% Srautų sąrašas ir aprašai generuojami iš Marijaus darbinio failo; jei
% keičiate diagramas, praneškite Marijui, kad srautų aprašai liktų suderinti.

\tikzset{
  dfdproc/.style={circle, draw, thick, align=center, font=\scriptsize,
    minimum size=2.4cm, text width=2.0cm, inner sep=0pt,
    execute at begin node={\hyphenpenalty=10000\exhyphenpenalty=10000}},
  dfdproc2/.style={dfdproc, minimum size=2.1cm, text width=1.75cm},
  dfdext/.style={rectangle, draw, thick, fill=gray!12, align=center,
    font=\scriptsize, minimum width=2.3cm, minimum height=0.9cm},
  dfdstore/.style={align=left, font=\scriptsize, text width=4.29cm,
    minimum height=0.5cm, inner xsep=3pt, inner ysep=0pt},
  dfdflow/.style={-{Stealth[length=2.2mm]}, semithick},
  srautas/.style={font=\scriptsize, align=center, inner sep=1.2pt},
  srautasv/.style={font=\scriptsize, inner sep=1.5pt},
}

\subsection{Duomenų srautų diagramos}\label{sec:dfd}

Duomenų srautų modelis sudarytas iš viršaus žemyn: kontekstinė diagrama
(0~lygis) rodo visą sistemą kaip vieną procesą, 1~lygio diagrama~– septynis
pagrindinius procesus ir devynias duomenų saugyklas (D1--D9), o 2~lygio diagramos
detalizuoja du sudėtingiausius procesus~– rezervacijas ir užsiėmimo vykdymą.
Kiekvienas srautas aprašytas \ref{sec:dfd-srautai}~poskyryje, o žemiausio
lygio procesų logika~– \ref{sec:dfd-pspec}~poskyryje.

Naudojamas Yourdono ir DeMarco žymėjimas: apskritimas~– procesas,
stačiakampis~– išorinis objektas, dvi lygiagrečios linijos~– duomenų
saugykla, rodyklė~– duomenų srautas.

Sistemos riba tokia pati kaip projekto plano funkcinių taškų (FP) skaičiavime:
sistemą sudaro centrinė dalis, žiniatinklio sąsaja ir vietinio grotuvo
programa. Todėl grotuvas yra sistemos viduje, o išoriniai objektai~– trys
naudotojų vaidmenys ir sporto erdvės įrenginiai, gaunantys vaizdą ir garsą:
ekranas ir garso įranga (radijo siųstuvas arba kolonėlės). Trenerio
mikrofonas diagramose nerodomas, nes trenerio balsas į radijo siųstuvą
patenka tiesiogiai, aplenkdamas sistemą. Budintis darbuotojas, kuriam pranešama apie
ryšio sutrikimą, atlieka administratoriaus vaidmenį.

\subsubsection{Kontekstinė diagrama}\label{sec:dfd-0}

\begin{figure}[H]
\centering
\begin{tikzpicture}[x=1cm,y=1cm]
  \node[dfdproc, minimum size=3.4cm, text width=2.8cm, font=\small] (S) at (0,0)
    {\textbf{0}\\[2pt]Užsiėmimų valdymo sistema};
  \node[dfdext] (T) at (0,4.6) {Treneris};
  \node[dfdext, minimum height=2.8cm] (D) at (-6.5,0) {Dalyvis};
  \node[dfdext, minimum height=2.8cm, minimum width=2.6cm] (A) at (6.5,0) {Administratorius};
  \node[dfdext] (E) at (-2.6,-4.6) {Ekranas};
  \node[dfdext, text width=2.6cm] (G) at (2.6,-4.6)
    {Garso įranga\\{\scriptsize (radijo siųstuvas arba kolonėlės)}};
  % Treneris
  \draw[dfdflow] (S.110 |- T.south) -- node[srautas, left]
    {I2 Trenerio duomenys\\ir komandos} (S.110);
  \draw[dfdflow] (S.70) -- node[srautas, right]
    {O3 Informacija\\treneriui} (S.70 |- T.south);
  % Dalyvis
  \draw[dfdflow] (D.east |- S.144) -- node[srautas, above]
    {I1 Dalyvio duomenys\\ir veiksmai} (S.144);
  \draw[dfdflow] (S.west) -- node[srautas, above]
    {O1 Informacija dalyviui} (D.east |- S.west);
  \draw[dfdflow] (S.216) -- node[srautas, above]
    {O2 Pranešimai dalyviui} (D.east |- S.216);
  % Administratorius
  \draw[dfdflow] (A.west |- S.36) -- node[srautas, above]
    {I3 Administravimo\\duomenys ir turinys} (S.36);
  \draw[dfdflow] (S.east) -- node[srautas, above]
    {O4 Administravimo\\informacija} (A.west |- S.east);
  \draw[dfdflow] (S.324) -- node[srautas, above]
    {O5 Pranešimas apie\\ryšio sutrikimą} (A.west |- S.324);
  % Įranga
  \draw[dfdflow] (S) -- node[srautas, left=3pt, pos=0.55]
    {O6 Pratimo vaizdas\\ir etapo informacija} (E);
  \draw[dfdflow] (S) -- node[srautas, right=3pt, pos=0.55]
    {O7 Garso\\instrukcijos} (G);
\end{tikzpicture}
\caption{Kontekstinė duomenų srautų diagrama (0 lygis)}\label{fig:dfd-0}
\end{figure}

Kontekstinėje diagramoje (\ref{fig:dfd-0}~pav.) parodyta dešimt srautų tarp
sistemos ir išorinių objektų. Tai apibendrinti srautai: 1~lygyje dauguma jų
išskaidomi į smulkesnius, pavyzdžiui, I1 „Dalyvio duomenys ir
veiksmai“~= I1.1 „Dalyvio paskyros duomenys“~+ I1.2 „Rezervacijos veiksmas“.
Sudėtis pateikta \ref{tab:dfd-0}~lentelėje: iš jos matyti, kad abiejų lygių
srautai subalansuoti.

\subsubsection{1 lygio diagrama}\label{sec:dfd-1}

Pirmojo lygio diagramoje (\ref{fig:dfd-1a} ir \ref{fig:dfd-1b}~pav.) sistema
išskaidyta į septynis procesus pagal sistemos vizijos dalis
(\ref{tab:dfd-procesai}~lentelė). Sistemos dalys pažymėtos projekto plano
įvertinimo kodais: P0--P7 dydis apskaičiuotas funkciniais taškais, o vietinio
grotuvo programa S2 įvertinta atskirai kaip FP nematoma programinė įranga. Duomenų saugyklos D1--D9 atitinka
funkcinių taškų skaičiavimo vidinius loginius failus
(\ref{tab:dfd-saugyklos}~lentelė), todėl abiejuose dokumentuose duomenys
sugrupuoti vienodai; saugyklų turinį detalizuoja ER diagrama (\ref{sec:er}~poskyris, \ref{fig:er}~pav.).

\begin{table}[H]
\centering
\caption{1 lygio procesai}\label{tab:dfd-procesai}
\small
\begin{tabular}{@{}l>{\raggedright\arraybackslash}p{6.2cm}>{\raggedright\arraybackslash}p{3.2cm}>{\raggedright\arraybackslash}p{3.6cm}@{}}
\toprule
\textbf{Nr.} & \textbf{Procesas} & \textbf{Sistemos dalis} & \textbf{Panaudojimo atvejai} \\
\midrule
1 & Tvarkyti paskyras, erdves ir įrangą & P0 & PA-01--PA-03 \\
2 & Tvarkyti pratimus, medžiagą ir programas & P1 & PA-04, PA-05 \\
3 & Planuoti užsiėmimus & P2 (planavimas) & PA-06--PA-08 \\
4 & Valdyti rezervacijas ir lankomumą & P3 & PA-09--PA-12, PA-18 \\
5 & Tvarkyti asmenines programas & P4 & PA-13, PA-14 \\
6 & Vykdyti užsiėmimą & P2 (vykdymas), P5, P6, S2 & PA-15, PA-16 \\
7 & Pateikti istoriją ir ataskaitas & P7 & PA-17 \\
\bottomrule
\end{tabular}
\end{table}

\begin{table}[H]
\centering
\caption{Duomenų saugyklos}\label{tab:dfd-saugyklos}
\small
\begin{tabular}{@{}l>{\raggedright\arraybackslash}p{3.6cm}l>{\raggedright\arraybackslash}p{8.4cm}@{}}
\toprule
\textbf{ID} & \textbf{Saugykla} & \textbf{FP failas} & \textbf{Turinys} \\
\midrule
% >>> GEN:stores
D1 & Naudotojų paskyros & P0.1 & Naudotojo paskyra, Dalyvio profilis, Trenerio profilis \\
D2 & Sporto erdvės & P0.2 & Sporto erdvė, Erdvės vaizdo ir garso profilis \\
D3 & Įranga & P0.3 & Įrenginys (su konfigūracijos pasu), Vaizdo įrenginio nustatymai, Garso įrenginio nustatymai \\
D4 & Pratimai & P1.1 & Pratimas \\
D5 & Mokomoji medžiaga & P1.2 & Mokomoji medžiaga (vaizdo ir garso įrašai) \\
D6 & Užsiėmimų programos & P1.3 & Užsiėmimo programa, Programos versija, Programos etapas, Etapo pratimas \\
D7 & Užsiėmimai & P2.1 & Užsiėmimas, Užsiėmimo nustatymai, Vykdymo žurnalas \\
D8 & Rezervacijos ir lankomumas & P3.1 & Rezervacija, Lankomumo įrašas \\
D9 & Asmeninės programos & P4.1 & Įvertinimas, Asmeninė programa, Programos pratimas, Pasiruošimo rekomendacija, Namų užduotis \\
D10 & Vietinė turinio kopija & -- & Grotuve saugoma užsiėmimo turinio kopija (tik 6 procese) \\
% <<< GEN:stores
\bottomrule
\end{tabular}
\end{table}

Kad diagrama būtų įskaitoma, ji padalyta į du paveikslus, o kai kurie
išoriniai objektai ir saugyklos pavaizduoti po kelis kartus~– taip rodyklės
nesikerta. Kelis kartus pavaizduoti simboliai pažymėti žvaigždute (*).
Srautai tarp procesų ir saugyklų diagramose pažymėti tik
identifikatoriais, o jų pavadinimai ir aprašai pateikti \ref{tab:dfd-v} ir
\ref{tab:dfd-l2}~lentelėse. Peržiūros užklausų parametrai (pavyzdžiui, laikotarpis ar filtras)
atskirais srautais nerodomi: jie yra atsakymo srauto dalis ir nurodyti jo
perdavimo sąlygose.

% >>> GEN:dfd-1a
\begin{figure}[p]
\centering
\begin{tikzpicture}[x=1cm,y=1cm]
  \node[dfdproc] (p1) at (0.3,0) {\textbf{1}\\Tvarkyti paskyras, erdves ir įrangą};
  \node[dfdext] (dal1) at (-6.85,2.4) {Dalyvis};
  \node[font=\scriptsize, inner sep=1.5pt, anchor=south east] at (dal1.south east) {*};
  \node[dfdext] (tre1) at (-6.85,0) {Treneris};
  \node[font=\scriptsize, inner sep=1.5pt, anchor=south east] at (tre1.south east) {*};
  \node[dfdext] (adm1) at (-6.85,-2.4) {Administratorius};
  \node[font=\scriptsize, inner sep=1.5pt, anchor=south east] at (adm1.south east) {*};
  \node[dfdstore] (d1) at (5.55,2) {\textbf{D1}~~Naudotojų paskyros\,*};
  \draw[thick] (d1.north west) -- (d1.north east) (d1.south west) -- (d1.south east);
  \node[dfdstore] (d2) at (5.55,1) {\textbf{D2}~~Sporto erdvės};
  \draw[thick] (d2.north west) -- (d2.north east) (d2.south west) -- (d2.south east);
  \node[dfdstore] (d3) at (5.55,0) {\textbf{D3}~~Įranga};
  \draw[thick] (d3.north west) -- (d3.north east) (d3.south west) -- (d3.south east);
  \node[dfdstore] (d7a) at (5.55,-1) {\textbf{D7}~~Užsiėmimai\,*};
  \draw[thick] (d7a.north west) -- (d7a.north east) (d7a.south west) -- (d7a.south east);
  \node[dfdstore] (d8a) at (5.55,-2) {\textbf{D8}~~Rezervacijos ir lankomumas\,*};
  \draw[thick] (d8a.north west) -- (d8a.north east) (d8a.south west) -- (d8a.south east);
  \node[dfdproc] (p2) at (0.3,-6.45) {\textbf{2}\\Tvarkyti pratimus, medžiagą ir programas};
  \node[dfdext] (tre2) at (-6.85,-5.25) {Treneris};
  \node[font=\scriptsize, inner sep=1.5pt, anchor=south east] at (tre2.south east) {*};
  \node[dfdext] (adm2) at (-6.85,-7.65) {Administratorius};
  \node[font=\scriptsize, inner sep=1.5pt, anchor=south east] at (adm2.south east) {*};
  \node[dfdstore] (d1c) at (5.55,-3.7) {\textbf{D1}~~Naudotojų paskyros\,*};
  \draw[thick] (d1c.north west) -- (d1c.north east) (d1c.south west) -- (d1c.south east);
  \node[dfdstore] (d7b) at (5.55,-4.8) {\textbf{D7}~~Užsiėmimai\,*};
  \draw[thick] (d7b.north west) -- (d7b.north east) (d7b.south west) -- (d7b.south east);
  \node[dfdstore] (d6) at (5.55,-5.9) {\textbf{D6}~~Užsiėmimų programos};
  \draw[thick] (d6.north west) -- (d6.north east) (d6.south west) -- (d6.south east);
  \node[dfdstore] (d5) at (5.55,-7) {\textbf{D5}~~Mokomoji medžiaga};
  \draw[thick] (d5.north west) -- (d5.north east) (d5.south west) -- (d5.south east);
  \node[dfdstore] (d9b) at (5.55,-8.1) {\textbf{D9}~~Asmeninės programos\,*};
  \draw[thick] (d9b.north west) -- (d9b.north east) (d9b.south west) -- (d9b.south east);
  \node[dfdstore] (d4) at (5.55,-9.2) {\textbf{D4}~~Pratimai\,*};
  \draw[thick] (d4.north west) -- (d4.north east) (d4.south west) -- (d4.south east);
  \node[dfdproc] (p5) at (0.3,-12.7) {\textbf{5}\\Tvarkyti asmenines programas};
  \node[dfdext] (tre5) at (-6.85,-11.5) {Treneris};
  \node[font=\scriptsize, inner sep=1.5pt, anchor=south east] at (tre5.south east) {*};
  \node[dfdext] (dal5) at (-6.85,-13.9) {Dalyvis};
  \node[font=\scriptsize, inner sep=1.5pt, anchor=south east] at (dal5.south east) {*};
  \node[dfdstore] (d4e) at (5.55,-10.7) {\textbf{D4}~~Pratimai\,*};
  \draw[thick] (d4e.north west) -- (d4e.north east) (d4e.south west) -- (d4e.south east);
  \node[dfdstore] (d9) at (5.55,-11.7) {\textbf{D9}~~Asmeninės programos\,*};
  \draw[thick] (d9.north west) -- (d9.north east) (d9.south west) -- (d9.south east);
  \node[dfdstore] (d7e) at (5.55,-12.7) {\textbf{D7}~~Užsiėmimai\,*};
  \draw[thick] (d7e.north west) -- (d7e.north east) (d7e.south west) -- (d7e.south east);
  \node[dfdstore] (d8e) at (5.55,-13.7) {\textbf{D8}~~Rezervacijos ir lankomumas\,*};
  \draw[thick] (d8e.north west) -- (d8e.north east) (d8e.south west) -- (d8e.south east);
  \node[dfdstore] (d1b) at (5.55,-14.7) {\textbf{D1}~~Naudotojų paskyros\,*};
  \draw[thick] (d1b.north west) -- (d1b.north east) (d1b.south west) -- (d1b.south east);
  \draw[dfdflow] (-5.7,2.25) -- node[srautas, pos=0.431, sloped, below] {I1.1 Dalyvio paskyros duomenys} (-0.7,0.66);
  \draw[dfdflow] (-0.42,0.96) -- node[srautas, pos=0.699, sloped, above] {O1.1 Dalyvio profilis} (-5.7,2.55);
  \draw[dfdflow] (-5.7,-0.15) -- node[srautas, pos=0.271, sloped, below] {I2.1 Trenerio\\paskyros duomenys} (-0.89,-0.15);
  \draw[dfdflow] (-0.89,0.15) -- node[srautas, pos=0.644, sloped, above] {O3.1 Trenerio profilis} (-5.7,0.15);
  \draw[dfdflow] (-5.7,-2.55) -- node[srautas, pos=0.373, sloped, below] {I3.1 Paskyros, erdvės ir įranga} (-0.42,-0.96);
  \draw[dfdflow] (-0.7,-0.66) -- node[srautas, pos=0.688, sloped, above] {O4.1 Paskyrų, erdvių\\ir įrangos sąrašai} (-5.7,-2.25);
  \draw[dfdflow] (1.28,0.69) -- node[srautasv, pos=0.759, sloped, below] {V1} (3.3,1.85);
  \draw[dfdflow] (3.3,2.15) -- node[srautasv, pos=0.136, sloped, above] {V2} (1.02,0.96);
  \draw[dfdflow] (1.49,0.14) -- node[srautasv, pos=0.726, sloped, below] {V3} (3.3,0.85);
  \draw[dfdflow] (3.3,1.15) -- node[srautasv, pos=0.417, sloped, above] {V4} (1.43,0.41);
  \draw[dfdflow] (1.43,-0.41) -- node[srautasv, pos=0.702, sloped, below] {V5} (3.3,-0.15);
  \draw[dfdflow] (3.3,0.15) -- node[srautasv, pos=0.492, sloped, above] {V6} (1.49,-0.14);
  \draw[dfdflow] (3.3,-1) -- node[srautasv, pos=0.172, sloped, below] {V7} (1.28,-0.69);
  \draw[dfdflow] (1.02,-0.96) -- node[srautasv, pos=0.777, sloped, below] {V8} (3.3,-2);
  \draw[dfdflow] (-5.7,-5.4) -- node[srautas, pos=0.367, sloped, below] {I2.2 Programos duomenys} (-0.83,-6.06);
  \draw[dfdflow] (-0.68,-5.76) -- node[srautas, pos=0.647, sloped, above] {O3.2 Pratimai,\\medžiaga ir programos} (-5.7,-5.1);
  \draw[dfdflow] (-5.7,-7.8) -- node[srautas, pos=0.226, sloped, below] {I3.2 Pratimai ir\\medžiagos failai} (-0.68,-7.14);
  \draw[dfdflow] (-0.83,-6.84) -- node[srautas, pos=0.636, sloped, above] {O4.2 Pratimai,\\medžiaga ir programos} (-5.7,-7.5);
  \draw[dfdflow] (1.02,-7.41) -- node[srautasv, pos=0.813, sloped, below] {V9} (3.3,-9.35);
  \draw[dfdflow] (3.3,-9.05) -- node[srautasv, pos=0.173, sloped, above] {V10} (1.26,-7.17);
  \draw[dfdflow] (1.48,-6.69) -- node[srautasv, pos=0.662, sloped, below] {V11} (3.3,-7.15);
  \draw[dfdflow] (3.3,-6.85) -- node[srautasv, pos=0.212, sloped, above] {V12} (1.5,-6.45);
  \draw[dfdflow] (1.48,-6.21) -- node[srautasv, pos=0.442, sloped, below] {V13} (3.3,-6.05);
  \draw[dfdflow] (3.3,-5.75) -- node[srautasv, pos=0.22, sloped, above] {V14} (1.4,-5.97);
  \draw[dfdflow] (3.3,-4.8) -- node[srautasv, pos=0.188, sloped, above] {V15} (1.26,-5.73);
  \draw[dfdflow] (3.3,-8.1) -- node[srautasv, pos=0.283, sloped, below] {V16} (1.4,-6.93);
  \draw[dfdflow] (3.3,-3.7) -- node[srautasv, pos=0.155, sloped, above] {V49} (1.02,-5.49);
  \draw[dfdflow] (-5.7,-11.65) -- node[srautas, pos=0.284, sloped, below] {I2.5 Įvertinimas ir\\asmeninė programa} (-0.83,-12.31);
  \draw[dfdflow] (-0.68,-12.01) -- node[srautas, pos=0.593, sloped, above] {O3.5 Asmeninės programos} (-5.7,-11.35);
  \draw[dfdflow] (-0.68,-13.39) -- node[srautas, pos=0.614, sloped, below] {O1.3 Asmeninė programa} (-5.7,-14.05);
  \draw[dfdflow] (-0.83,-13.09) -- node[srautas, pos=0.65, sloped, above] {O2.3 Programos\\pakeitimo pranešimas} (-5.7,-13.75);
  \draw[dfdflow] (1.46,-12.4) -- node[srautasv, pos=0.672, sloped, below] {V17} (3.3,-11.85);
  \draw[dfdflow] (3.3,-11.55) -- node[srautasv, pos=0.207, sloped, above] {V18} (1.34,-12.1);
  \draw[dfdflow] (3.3,-10.7) -- node[srautasv, pos=0.17, sloped, above] {V19} (1.09,-11.8);
  \draw[dfdflow] (3.3,-14.7) -- node[srautasv, pos=0.17, sloped, below] {V20} (1.09,-13.6);
  \draw[dfdflow] (3.3,-12.7) -- node[srautasv, pos=0.233, sloped, below] {V50} (1.5,-12.7);
  \draw[dfdflow] (3.3,-13.7) -- node[srautasv, pos=0.214, sloped, below] {V51} (1.41,-13.15);
\end{tikzpicture}
\caption{1 lygio duomenų srautų diagrama (1, 2 ir 5 procesai)}\label{fig:dfd-1a}
\end{figure}
% <<< GEN:dfd-1a

% >>> GEN:dfd-1b
\begin{figure}[p]
\centering
\begin{tikzpicture}[x=1cm,y=1cm]
  \node[dfdproc] (p3) at (0.3,0) {\textbf{3}\\Planuoti užsiėmimus};
  \node[dfdext] (tre3) at (-6.85,2.4) {Treneris};
  \node[font=\scriptsize, inner sep=1.5pt, anchor=south east] at (tre3.south east) {*};
  \node[dfdext] (adm3) at (-6.85,0) {Administratorius};
  \node[font=\scriptsize, inner sep=1.5pt, anchor=south east] at (adm3.south east) {*};
  \node[dfdext] (dal3) at (-6.85,-2.4) {Dalyvis};
  \node[font=\scriptsize, inner sep=1.5pt, anchor=south east] at (dal3.south east) {*};
  \node[dfdstore] (d1) at (5.55,2) {\textbf{D1}~~Naudotojų paskyros\,*};
  \draw[thick] (d1.north west) -- (d1.north east) (d1.south west) -- (d1.south east);
  \node[dfdstore] (d2) at (5.55,1) {\textbf{D2}~~Sporto erdvės\,*};
  \draw[thick] (d2.north west) -- (d2.north east) (d2.south west) -- (d2.south east);
  \node[dfdstore] (d6) at (5.55,0) {\textbf{D6}~~Užsiėmimų programos\,*};
  \draw[thick] (d6.north west) -- (d6.north east) (d6.south west) -- (d6.south east);
  \node[dfdstore] (d7) at (5.55,-1) {\textbf{D7}~~Užsiėmimai\,*};
  \draw[thick] (d7.north west) -- (d7.north east) (d7.south west) -- (d7.south east);
  \node[dfdstore] (d8) at (5.55,-2) {\textbf{D8}~~Rezervacijos ir lankomumas\,*};
  \draw[thick] (d8.north west) -- (d8.north east) (d8.south west) -- (d8.south east);
  \node[dfdproc] (p4) at (0.3,-6.8) {\textbf{4}\\Valdyti rezervacijas ir lankomumą};
  \node[dfdext, minimum height=1.3cm] (dal4) at (-6.85,-4.4) {Dalyvis};
  \node[font=\scriptsize, inner sep=1.5pt, anchor=south east] at (dal4.south east) {*};
  \node[dfdext] (tre4) at (-6.85,-6.8) {Treneris};
  \node[font=\scriptsize, inner sep=1.5pt, anchor=south east] at (tre4.south east) {*};
  \node[dfdext] (adm4) at (-6.85,-9.2) {Administratorius};
  \node[font=\scriptsize, inner sep=1.5pt, anchor=south east] at (adm4.south east) {*};
  \node[dfdstore] (d7b) at (5.55,-5.3) {\textbf{D7}~~Užsiėmimai\,*};
  \draw[thick] (d7b.north west) -- (d7b.north east) (d7b.south west) -- (d7b.south east);
  \node[dfdstore] (d8b) at (5.55,-6.3) {\textbf{D8}~~Rezervacijos ir lankomumas\,*};
  \draw[thick] (d8b.north west) -- (d8b.north east) (d8b.south west) -- (d8b.south east);
  \node[dfdstore] (d1b) at (5.55,-7.3) {\textbf{D1}~~Naudotojų paskyros\,*};
  \draw[thick] (d1b.north west) -- (d1b.north east) (d1b.south west) -- (d1b.south east);
  \node[dfdstore] (d2b) at (5.55,-8.3) {\textbf{D2}~~Sporto erdvės\,*};
  \draw[thick] (d2b.north west) -- (d2b.north east) (d2b.south west) -- (d2b.south east);
  \node[dfdproc] (p6) at (0.3,-13.1) {\textbf{6}\\Vykdyti užsiėmimą};
  \node[dfdext] (tre6) at (-6.85,-10.7) {Treneris};
  \node[font=\scriptsize, inner sep=1.5pt, anchor=south east] at (tre6.south east) {*};
  \node[dfdext] (adm6) at (-6.85,-12.3) {Administratorius};
  \node[font=\scriptsize, inner sep=1.5pt, anchor=south east] at (adm6.south east) {*};
  \node[dfdext] (ekr6) at (-6.85,-13.9) {Ekranas};
  \node[dfdext] (gar6) at (-6.85,-15.5) {Garso įranga};
  \node[dfdstore] (d7c) at (5.55,-10.6) {\textbf{D7}~~Užsiėmimai\,*};
  \draw[thick] (d7c.north west) -- (d7c.north east) (d7c.south west) -- (d7c.south east);
  \node[dfdstore] (d6c) at (5.55,-11.6) {\textbf{D6}~~Užsiėmimų programos\,*};
  \draw[thick] (d6c.north west) -- (d6c.north east) (d6c.south west) -- (d6c.south east);
  \node[dfdstore] (d5c) at (5.55,-12.6) {\textbf{D5}~~Mokomoji medžiaga};
  \draw[thick] (d5c.north west) -- (d5c.north east) (d5c.south west) -- (d5c.south east);
  \node[dfdstore] (d4c) at (5.55,-13.6) {\textbf{D4}~~Pratimai};
  \draw[thick] (d4c.north west) -- (d4c.north east) (d4c.south west) -- (d4c.south east);
  \node[dfdstore] (d3c) at (5.55,-14.6) {\textbf{D3}~~Įranga};
  \draw[thick] (d3c.north west) -- (d3c.north east) (d3c.south west) -- (d3c.south east);
  \node[dfdstore] (d2c) at (5.55,-15.6) {\textbf{D2}~~Sporto erdvės\,*};
  \draw[thick] (d2c.north west) -- (d2c.north east) (d2c.south west) -- (d2c.south east);
  \node[dfdproc] (p7) at (0.3,-18.2) {\textbf{7}\\Pateikti istoriją ir ataskaitas};
  \node[dfdext] (tre7) at (-6.85,-17) {Treneris};
  \node[font=\scriptsize, inner sep=1.5pt, anchor=south east] at (tre7.south east) {*};
  \node[dfdext] (adm7) at (-6.85,-19.4) {Administratorius};
  \node[font=\scriptsize, inner sep=1.5pt, anchor=south east] at (adm7.south east) {*};
  \node[dfdstore] (d7d) at (5.55,-17.2) {\textbf{D7}~~Užsiėmimai\,*};
  \draw[thick] (d7d.north west) -- (d7d.north east) (d7d.south west) -- (d7d.south east);
  \node[dfdstore] (d8d) at (5.55,-18.2) {\textbf{D8}~~Rezervacijos ir lankomumas\,*};
  \draw[thick] (d8d.north west) -- (d8d.north east) (d8d.south west) -- (d8d.south east);
  \node[dfdstore] (d2d) at (5.55,-19.2) {\textbf{D2}~~Sporto erdvės\,*};
  \draw[thick] (d2d.north west) -- (d2d.north east) (d2d.south west) -- (d2d.south east);
  \draw[dfdflow] (-5.7,2.25) -- node[srautas, pos=0.356, sloped, below] {I2.3 Užsiėmimo duomenys} (-0.7,0.66);
  \draw[dfdflow] (-0.42,0.96) -- node[srautas, pos=0.66, sloped, above] {O3.3 Užsiėmimų sąrašas} (-5.7,2.55);
  \draw[dfdflow] (-5.7,-0.15) -- node[srautas, pos=0.37, sloped, below] {I3.3 Užsiėmimo duomenys} (-0.89,-0.15);
  \draw[dfdflow] (-0.89,0.15) -- node[srautas, pos=0.61, sloped, above] {O4.3 Užsiėmimų sąrašas} (-5.7,0.15);
  \draw[dfdflow] (-0.42,-0.96) -- node[srautas, pos=0.686, sloped, above] {O2.1 Užsiėmimo\\pakeitimo pranešimas} (-5.7,-2.4);
  \draw[dfdflow] (1.44,-0.36) -- node[srautasv, pos=0.702, sloped, below] {V21} (3.3,-1.15);
  \draw[dfdflow] (3.3,-0.85) -- node[srautasv, pos=0.229, sloped, above] {V22} (1.5,-0.06);
  \draw[dfdflow] (3.3,2) -- node[srautasv, pos=0.17, sloped, above] {V23} (1.09,0.9);
  \draw[dfdflow] (3.3,1) -- node[srautasv, pos=0.211, sloped, above] {V24} (1.37,0.54);
  \draw[dfdflow] (3.3,0) -- node[srautasv, pos=0.245, sloped, above] {V25} (1.48,0.24);
  \draw[dfdflow] (3.3,-2.15) -- node[srautasv, pos=0.163, sloped, below] {V26} (1.02,-0.96);
  \draw[dfdflow] (1.3,-0.66) -- node[srautasv, pos=0.311, sloped, above] {V27} (3.3,-1.85);
  \draw[dfdflow] (-5.7,-4.84) -- node[srautas, pos=0.366, sloped, below] {I1.2 Rezervacijos veiksmas} (-0.84,-6.44);
  \draw[dfdflow] (-0.7,-6.14) -- node[srautas, pos=0.551, sloped, above] {O1.2 Tvarkaraštis ir rezervacijos} (-5.7,-4.4);
  \draw[dfdflow] (-0.42,-5.84) -- node[srautas, pos=0.574, sloped, above] {O2.2 Priminimas apie užsiėmimą} (-5.7,-3.96);
  \draw[dfdflow] (-5.7,-6.95) -- node[srautas, pos=0.244, sloped, below] {I2.4 Priskyrimas\\ir lankomumas} (-0.88,-7.04);
  \draw[dfdflow] (-0.9,-6.74) -- node[srautas, pos=0.653, sloped, above] {O3.4 Dalyvių sąrašai} (-5.7,-6.65);
  \draw[dfdflow] (-5.7,-9.35) -- node[srautas, pos=0.308, sloped, below] {I3.4 Dalyvių priskyrimas} (-0.42,-7.76);
  \draw[dfdflow] (-0.7,-7.46) -- node[srautas, pos=0.596, sloped, above] {O4.4 Užsiregistravę dalyviai} (-5.7,-9.05);
  \draw[dfdflow] (1.5,-6.72) -- node[srautasv, pos=0.651, sloped, below] {V28} (3.3,-6.45);
  \draw[dfdflow] (3.3,-6.15) -- node[srautasv, pos=0.224, sloped, above] {V29} (1.44,-6.42);
  \draw[dfdflow] (3.3,-5.3) -- node[srautasv, pos=0.193, sloped, above] {V30} (1.29,-6.12);
  \draw[dfdflow] (3.3,-7.3) -- node[srautasv, pos=0.228, sloped, below] {V31} (1.48,-7.02);
  \draw[dfdflow] (3.3,-8.3) -- node[srautasv, pos=0.193, sloped, below] {V32} (1.29,-7.48);
  \draw[dfdflow] (-5.7,-10.85) -- node[srautas, pos=0.342, sloped, below] {I2.6 Vykdymo komandos} (-0.7,-12.44);
  \draw[dfdflow] (-0.42,-12.14) -- node[srautas, pos=0.582, sloped, above] {O3.6 Valdymo pulto duomenys} (-5.7,-10.55);
  \draw[dfdflow] (-0.84,-12.74) -- node[srautas, pos=0.671, sloped, below] {O5 Pranešimas apie\\ryšio sutrikimą} (-5.7,-12.3);
  \draw[dfdflow] (-0.84,-13.46) -- node[srautas, pos=0.67, sloped, below] {O6 Pratimo vaizdas\\ir etapo informacija} (-5.7,-13.9);
  \draw[dfdflow] (-0.42,-14.06) -- node[srautas, pos=0.685, sloped, above] {O7 Garso instrukcijos} (-5.7,-15.5);
  \draw[dfdflow] (3.3,-10.75) -- node[srautasv, pos=0.183, sloped, below] {V33} (1.3,-12.44);
  \draw[dfdflow] (1.02,-12.14) -- node[srautasv, pos=0.777, sloped, above] {V34} (3.3,-10.45);
  \draw[dfdflow] (3.3,-11.6) -- node[srautasv, pos=0.249, sloped, below] {V35} (1.44,-12.74);
  \draw[dfdflow] (3.3,-12.6) -- node[srautasv, pos=0.26, sloped, below] {V36} (1.5,-13.04);
  \draw[dfdflow] (3.3,-14.6) -- node[srautasv, pos=0.194, sloped, below] {V37} (1.3,-13.76);
  \draw[dfdflow] (3.3,-13.6) -- node[srautasv, pos=0.228, sloped, below] {V47} (1.48,-13.34);
  \draw[dfdflow] (3.3,-15.6) -- node[srautasv, pos=0.153, sloped, below] {V48} (1.02,-14.06);
  \draw[dfdflow] (-0.77,-17.66) -- node[srautas, pos=0.604, sloped, above] {O3.7 Istorija ir ataskaitos} (-5.7,-17);
  \draw[dfdflow] (-0.77,-18.74) -- node[srautas, pos=0.604, sloped, above] {O4.5 Istorija ir ataskaitos} (-5.7,-19.4);
  \draw[dfdflow] (3.3,-17.2) -- node[srautasv, pos=0.214, sloped, above] {V38} (1.41,-17.75);
  \draw[dfdflow] (3.3,-18.2) -- node[srautasv, pos=0.233, sloped, below] {V39} (1.5,-18.2);
  \draw[dfdflow] (3.3,-19.2) -- node[srautasv, pos=0.214, sloped, below] {V40} (1.41,-18.65);
\end{tikzpicture}
\caption{1 lygio duomenų srautų diagrama (3, 4, 6 ir 7 procesai)}\label{fig:dfd-1b}
\end{figure}
% <<< GEN:dfd-1b

\clearpage
\subsubsection{2 lygio diagramos}\label{sec:dfd-2}

Detalizuoti du procesai, kuriuose daugiausia taisyklių: 4~procesas
„Valdyti rezervacijas ir lankomumą“ ir 6~procesas „Vykdyti užsiėmimą“.
Rezervacijų ir lankomumo procesas (\ref{fig:dfd-2-4}~pav.) išskaidytas pagal dalyvių
valdymo funkcijas: jame tikrinamos laisvos vietos, užsiėmimo pradžia ir
paskyros būsena, todėl rezervavimui sudaryta ir sprendimų lentelė. Užsiėmimo
vykdymo procesas (\ref{fig:dfd-2-6}~pav.) vyksta realiuoju laiku. Turinys
iš anksto atsisiunčiamas į grotuvą ir saugomas vietinėje kopijoje D10, todėl nutrūkus ryšiui su serveriu grotuvas užbaigia pradėtą
etapą, o serveris, nebegaudamas grotuvo signalo, praneša budinčiam
darbuotojui. D10 naudojama tik šiame procese, todėl 1~lygio diagramoje jos
nėra.

Antrajame lygyje dalis 1~lygio srautų išskaidyta (pavyzdžiui, V29 į
V29.1--V29.7), o nauji 6~proceso vidiniai srautai (tarp jo dalių ir saugyklos
D10) pažymėti V41--V46. Neskaidomi 1~lygio srautai 2~lygio diagramose rodomi
tuo pačiu identifikatoriumi ir aprašyti \ref{tab:dfd-io} ir
\ref{tab:dfd-v}~lentelėse.

% >>> GEN:dfd-2-4
\begin{landscape}
\begin{figure}[p]
\centering
\begin{tikzpicture}[x=1cm,y=1cm]
  \node[dfdproc2] (p41) at (-4.9,5) {\textbf{4.1}\\Pateikti tvarkaraštį};
  \node[dfdext] (dal41) at (-10.9,5) {Dalyvis};
  \node[font=\scriptsize, inner sep=1.5pt, anchor=south east] at (dal41.south east) {*};
  \node[dfdproc2] (p46) at (4.9,5) {\textbf{4.6}\\Pateikti rezervacijas ir dalyvių sąrašus};
  \node[dfdext] (dal46) at (10.9,6.3) {Dalyvis};
  \node[font=\scriptsize, inner sep=1.5pt, anchor=south east] at (dal46.south east) {*};
  \node[dfdext] (tre46) at (10.9,5) {Treneris};
  \node[font=\scriptsize, inner sep=1.5pt, anchor=south east] at (tre46.south east) {*};
  \node[dfdext] (adm46) at (10.9,3.7) {Administratorius};
  \node[font=\scriptsize, inner sep=1.5pt, anchor=south east] at (adm46.south east) {*};
  \node[dfdstore] (a7) at (0,6.2) {\textbf{D7}~~Užsiėmimai\,*};
  \draw[thick] (a7.north west) -- (a7.north east) (a7.south west) -- (a7.south east);
  \node[dfdstore] (a8) at (0,5.4) {\textbf{D8}~~Rezervacijos ir lankomumas\,*};
  \draw[thick] (a8.north west) -- (a8.north east) (a8.south west) -- (a8.south east);
  \node[dfdstore] (a1) at (0,4.6) {\textbf{D1}~~Naudotojų paskyros\,*};
  \draw[thick] (a1.north west) -- (a1.north east) (a1.south west) -- (a1.south east);
  \node[dfdstore] (a2) at (0,3.8) {\textbf{D2}~~Sporto erdvės\,*};
  \draw[thick] (a2.north west) -- (a2.north east) (a2.south west) -- (a2.south east);
  \node[dfdproc2] (p42) at (-4.9,1) {\textbf{4.2}\\Rezervuoti vietą};
  \node[dfdext] (dal42) at (-10.9,1) {Dalyvis};
  \node[font=\scriptsize, inner sep=1.5pt, anchor=south east] at (dal42.south east) {*};
  \node[dfdproc2] (p44) at (4.9,1) {\textbf{4.4}\\Priskirti ar pašalinti dalyvį};
  \node[dfdext] (tre44) at (10.9,2.3) {Treneris};
  \node[font=\scriptsize, inner sep=1.5pt, anchor=south east] at (tre44.south east) {*};
  \node[dfdext] (adm44) at (10.9,-0.3) {Administratorius};
  \node[font=\scriptsize, inner sep=1.5pt, anchor=south east] at (adm44.south east) {*};
  \node[dfdstore] (b7) at (0,1.8) {\textbf{D7}~~Užsiėmimai\,*};
  \draw[thick] (b7.north west) -- (b7.north east) (b7.south west) -- (b7.south east);
  \node[dfdstore] (b8) at (0,1) {\textbf{D8}~~Rezervacijos ir lankomumas\,*};
  \draw[thick] (b8.north west) -- (b8.north east) (b8.south west) -- (b8.south east);
  \node[dfdstore] (b1) at (0,0.2) {\textbf{D1}~~Naudotojų paskyros\,*};
  \draw[thick] (b1.north west) -- (b1.north east) (b1.south west) -- (b1.south east);
  \node[dfdproc2] (p43) at (-4.9,-2.6) {\textbf{4.3}\\Atšaukti rezervaciją};
  \node[dfdext] (dal43) at (-10.9,-2.6) {Dalyvis};
  \node[font=\scriptsize, inner sep=1.5pt, anchor=south east] at (dal43.south east) {*};
  \node[dfdproc2] (p45) at (4.9,-2.6) {\textbf{4.5}\\Žymėti lankomumą};
  \node[dfdext] (tre45) at (10.9,-2.6) {Treneris};
  \node[font=\scriptsize, inner sep=1.5pt, anchor=south east] at (tre45.south east) {*};
  \node[dfdstore] (c7) at (0,-2.2) {\textbf{D7}~~Užsiėmimai\,*};
  \draw[thick] (c7.north west) -- (c7.north east) (c7.south west) -- (c7.south east);
  \node[dfdstore] (c8) at (0,-3) {\textbf{D8}~~Rezervacijos ir lankomumas\,*};
  \draw[thick] (c8.north west) -- (c8.north east) (c8.south west) -- (c8.south east);
  \node[dfdproc2] (p47) at (-4.9,-5.8) {\textbf{4.7}\\Siųsti priminimus};
  \node[dfdext] (dal47) at (-10.9,-5.8) {Dalyvis};
  \node[font=\scriptsize, inner sep=1.5pt, anchor=south east] at (dal47.south east) {*};
  \node[dfdstore] (e7) at (0,-4.6) {\textbf{D7}~~Užsiėmimai\,*};
  \draw[thick] (e7.north west) -- (e7.north east) (e7.south west) -- (e7.south east);
  \node[dfdstore] (e8) at (0,-5.4) {\textbf{D8}~~Rezervacijos ir lankomumas\,*};
  \draw[thick] (e8.north west) -- (e8.north east) (e8.south west) -- (e8.south east);
  \node[dfdstore] (e1) at (0,-6.2) {\textbf{D1}~~Naudotojų paskyros\,*};
  \draw[thick] (e1.north west) -- (e1.north east) (e1.south west) -- (e1.south east);
  \node[dfdstore] (e2) at (0,-7) {\textbf{D2}~~Sporto erdvės\,*};
  \draw[thick] (e2.north west) -- (e2.north east) (e2.south west) -- (e2.south east);
  \draw[dfdflow] (-5.95,5) -- node[srautas, pos=0.577, sloped, above] {O1.2.1 Tvarkaraštis} (-9.75,5);
  \draw[dfdflow] (-2.25,5.4) -- node[srautasv, pos=0.325, sloped, above] {V29.1} (-3.87,5.18);
  \draw[dfdflow] (-2.25,6.2) -- node[srautasv, pos=0.283, sloped, above] {V30.1} (-4,5.54);
  \draw[dfdflow] (-2.25,4.6) -- node[srautasv, pos=0.325, sloped, below] {V31.1} (-3.87,4.82);
  \draw[dfdflow] (-2.25,3.8) -- node[srautasv, pos=0.283, sloped, below] {V32.1} (-4,4.46);
  \draw[dfdflow] (5.77,5.58) -- node[srautas, pos=0.619, sloped, above] {O1.2.4 Rezervacijų\\istorija} (9.75,6.3);
  \draw[dfdflow] (5.95,5) -- node[srautas, pos=0.532, sloped, above] {O3.4.3 Dalyvių sąrašai} (9.75,5);
  \draw[dfdflow] (5.77,4.42) -- node[srautas, pos=0.582, sloped, above] {O4.4.2 Užsiregistravę\\dalyviai} (9.75,3.7);
  \draw[dfdflow] (2.25,5.4) -- node[srautasv, pos=0.325, sloped, above] {V29.6} (3.87,5.18);
  \draw[dfdflow] (2.25,6.2) -- node[srautasv, pos=0.283, sloped, above] {V30.6} (4,5.54);
  \draw[dfdflow] (2.25,4.6) -- node[srautasv, pos=0.325, sloped, below] {V31.4} (3.87,4.82);
  \draw[dfdflow] (2.25,3.8) -- node[srautasv, pos=0.283, sloped, below] {V32.2} (4,4.46);
  \draw[dfdflow] (-9.75,0.85) -- node[srautas, pos=0.354, sloped, below] {I1.2.1 Rezervavimo\\užklausa} (-5.94,0.85);
  \draw[dfdflow] (-5.94,1.15) -- node[srautas, pos=0.575, sloped, above] {O1.2.2 Rezervavimo\\rezultatas} (-9.75,1.15);
  \draw[dfdflow] (-3.86,0.85) -- node[srautasv, pos=0.557, sloped, below] {V28.1} (-2.25,0.85);
  \draw[dfdflow] (-2.25,1.15) -- node[srautasv, pos=0.33, sloped, above] {V29.2} (-3.86,1.15);
  \draw[dfdflow] (-2.25,1.8) -- node[srautasv, pos=0.306, sloped, above] {V30.2} (-3.95,1.45);
  \draw[dfdflow] (-2.25,0.2) -- node[srautasv, pos=0.306, sloped, below] {V31.2} (-3.95,0.55);
  \draw[dfdflow] (9.75,2.15) -- node[srautas, pos=0.276, sloped, below] {I2.4.1 Dalyvio\\priskyrimas} (5.86,1.44);
  \draw[dfdflow] (9.75,-0.45) -- node[srautas, pos=0.408, sloped, below] {I3.4 Dalyvių priskyrimas} (5.65,0.26);
  \draw[dfdflow] (5.65,1.73) -- node[srautas, pos=0.639, sloped, above] {O3.4.1 Priskyrimo\\rezultatas} (9.75,2.45);
  \draw[dfdflow] (5.86,0.56) -- node[srautas, pos=0.62, sloped, above] {O4.4.1 Priskyrimo\\rezultatas} (9.75,-0.15);
  \draw[dfdflow] (3.86,0.85) -- node[srautasv, pos=0.557, sloped, below] {V28.3} (2.25,0.85);
  \draw[dfdflow] (2.25,1.15) -- node[srautasv, pos=0.33, sloped, above] {V29.4} (3.86,1.15);
  \draw[dfdflow] (2.25,1.8) -- node[srautasv, pos=0.306, sloped, above] {V30.4} (3.95,1.45);
  \draw[dfdflow] (2.25,0.2) -- node[srautasv, pos=0.306, sloped, below] {V31.3} (3.95,0.55);
  \draw[dfdflow] (-9.75,-2.75) -- node[srautas, pos=0.321, sloped, below] {I1.2.2 Atšaukimo\\užklausa} (-5.94,-2.75);
  \draw[dfdflow] (-5.94,-2.45) -- node[srautas, pos=0.608, sloped, above] {O1.2.3 Atšaukimo\\rezultatas} (-9.75,-2.45);
  \draw[dfdflow] (-3.91,-2.96) -- node[srautasv, pos=0.555, sloped, below] {V28.2} (-2.25,-3.15);
  \draw[dfdflow] (-2.25,-2.85) -- node[srautasv, pos=0.349, sloped, above] {V29.3} (-3.85,-2.66);
  \draw[dfdflow] (-2.25,-2.2) -- node[srautasv, pos=0.324, sloped, above] {V30.3} (-3.88,-2.36);
  \draw[dfdflow] (9.75,-2.75) -- node[srautas, pos=0.346, sloped, below] {I2.4.2 Lankomumo\\žymos} (5.94,-2.75);
  \draw[dfdflow] (5.94,-2.45) -- node[srautas, pos=0.54, sloped, above] {O3.4.2 Lankomumo\\išsaugojimo rezultatas} (9.75,-2.45);
  \draw[dfdflow] (3.91,-2.96) -- node[srautasv, pos=0.555, sloped, below] {V28.4} (2.25,-3.15);
  \draw[dfdflow] (2.25,-2.85) -- node[srautasv, pos=0.349, sloped, above] {V29.5} (3.85,-2.66);
  \draw[dfdflow] (2.25,-2.2) -- node[srautasv, pos=0.324, sloped, above] {V30.5} (3.88,-2.36);
  \draw[dfdflow] (-5.95,-5.8) -- node[srautas, pos=0.625, sloped, above] {O2.2 Priminimas\\apie užsiėmimą} (-9.75,-5.8);
  \draw[dfdflow] (-3.85,-5.76) -- node[srautasv, pos=0.557, sloped, below] {V28.5} (-2.25,-5.55);
  \draw[dfdflow] (-2.25,-5.25) -- node[srautasv, pos=0.318, sloped, above] {V29.7} (-3.91,-5.46);
  \draw[dfdflow] (-2.25,-4.6) -- node[srautasv, pos=0.279, sloped, above] {V30.7} (-4.06,-5.16);
  \draw[dfdflow] (-2.25,-6.2) -- node[srautasv, pos=0.324, sloped, below] {V31.5} (-3.88,-6.06);
  \draw[dfdflow] (-2.25,-7) -- node[srautasv, pos=0.283, sloped, below] {V32.3} (-4.01,-6.36);
\end{tikzpicture}
\caption{2 lygio duomenų srautų diagrama: 4 procesas „Valdyti rezervacijas ir lankomumą“}\label{fig:dfd-2-4}
\end{figure}
\end{landscape}
% <<< GEN:dfd-2-4

% >>> GEN:dfd-2-6
\begin{figure}[p]
\centering
\begin{tikzpicture}[x=1cm,y=1cm]
  \node[dfdstore] (d6) at (-2.6,7.6) {\textbf{D6}~~Užsiėmimų programos};
  \draw[thick] (d6.north west) -- (d6.north east) (d6.south west) -- (d6.south east);
  \node[dfdstore] (d5) at (-2.6,6.8) {\textbf{D5}~~Mokomoji medžiaga};
  \draw[thick] (d5.north west) -- (d5.north east) (d5.south west) -- (d5.south east);
  \node[dfdstore] (d4) at (-2.6,6) {\textbf{D4}~~Pratimai};
  \draw[thick] (d4.north west) -- (d4.north east) (d4.south west) -- (d4.south east);
  \node[dfdstore] (d3) at (-2.6,5.2) {\textbf{D3}~~Įranga};
  \draw[thick] (d3.north west) -- (d3.north east) (d3.south west) -- (d3.south east);
  \node[dfdstore] (d7a) at (-2.6,4.4) {\textbf{D7}~~Užsiėmimai\,*};
  \draw[thick] (d7a.north west) -- (d7a.north east) (d7a.south west) -- (d7a.south east);
  \node[dfdproc2] (p61) at (2.6,6) {\textbf{6.1}\\Parengti turinį grotuvui};
  \node[dfdstore] (d10) at (2.6,3.4) {\textbf{D10}~~Vietinė turinio kopija};
  \draw[thick] (d10.north west) -- (d10.north east) (d10.south west) -- (d10.south east);
  \node[dfdext] (tre) at (-7,2) {Treneris};
  \node[font=\scriptsize, inner sep=1.5pt, anchor=south east] at (tre.south east) {*};
  \node[dfdproc2] (p62) at (-1.4,2) {\textbf{6.2}\\Priimti vykdymo komandas};
  \node[dfdstore] (d7) at (-1.4,-0.6) {\textbf{D7}~~Užsiėmimai\,*};
  \draw[thick] (d7.north west) -- (d7.north east) (d7.south west) -- (d7.south east);
  \node[dfdproc2] (p63) at (3.6,-0.6) {\textbf{6.3}\\Vykdyti programos eigą};
  \node[dfdproc2] (p66) at (-1.4,-3.6) {\textbf{6.6}\\Stebėti eigą ir ryšį su grotuvu};
  \node[dfdext] (tre2) at (-7,-2.7) {Treneris};
  \node[font=\scriptsize, inner sep=1.5pt, anchor=south east] at (tre2.south east) {*};
  \node[dfdext] (adm2) at (-7,-4.5) {Administratorius};
  \node[dfdstore] (d2) at (-3.6,-6.4) {\textbf{D2}~~Sporto erdvės};
  \draw[thick] (d2.north west) -- (d2.north east) (d2.south west) -- (d2.south east);
  \node[dfdproc2] (p64) at (1,-6) {\textbf{6.4}\\Rodyti etapą ekrane};
  \node[dfdproc2] (p65) at (5.4,-6) {\textbf{6.5}\\Perduoti garso instrukcijas};
  \node[dfdext] (ekr) at (1,-8.6) {Ekranas};
  \node[dfdext] (gar) at (5.4,-8.6) {Garso įranga};
  \draw[dfdflow] (-0.35,4.25) -- node[srautasv, pos=0.213, sloped, below] {V33.1} (1.97,5.16);
  \draw[dfdflow] (-0.35,7.6) -- node[srautasv, pos=0.178, sloped, above] {V35} (1.84,6.72);
  \draw[dfdflow] (-0.35,6.8) -- node[srautasv, pos=0.208, sloped, above] {V36} (1.61,6.36);
  \draw[dfdflow] (-0.35,5.2) -- node[srautasv, pos=0.179, sloped, below] {V37} (1.58,5.76);
  \draw[dfdflow] (-0.35,6) -- node[srautasv, pos=0.221, sloped, below] {V47} (1.55,6.06);
  \draw[dfdflow] (1.7,5.46) -- node[srautasv, pos=0.399, sloped, above] {V34.3} (-0.35,4.55);
  \draw[dfdflow] (2.6,4.95) -- node[srautasv, pos=0.606, anchor=west] {V41} (2.6,3.65);
  \draw[dfdflow] (2.6,3.15) -- node[srautasv, pos=0.101, anchor=west] {V42} (3.15,0.35);
  \draw[dfdflow] (-5.85,1.85) -- node[srautas, pos=0.304, sloped, below] {I2.6 Vykdymo\\komandos} (-2.44,1.85);
  \draw[dfdflow] (-2.44,2.15) -- node[srautas, pos=0.567, sloped, above] {O3.6.1 Komandos\\rezultatas} (-5.85,2.15);
  \draw[dfdflow] (-1.55,-0.35) -- node[srautasv, pos=0.224, anchor=east] {V33.2} (-1.55,0.96);
  \draw[dfdflow] (-1.25,0.96) -- node[srautasv, pos=0.615, anchor=west] {V34.1} (-1.25,-0.35);
  \draw[dfdflow] (-0.47,1.52) -- node[srautas, pos=0.47, sloped, above] {V43 Vykdymo komanda} (2.67,-0.12);
  \draw[dfdflow] (3.14,-1.55) -- node[srautas, pos=0.635, anchor=west] {V44 Etapo vaizdas} (1.46,-5.05);
  \draw[dfdflow] (3.93,-1.6) -- node[srautas, pos=0.496, anchor=west] {V45 Etapo garsas} (5.07,-5);
  \draw[dfdflow] (2.7,-1.14) -- node[srautas, pos=0.5, sloped, above] {V46 Grotuvo būsena} (-0.5,-3.06);
  \draw[dfdflow] (1,-7.05) -- node[srautas, pos=0.416, anchor=west] {O6 Pratimo vaizdas\\ir etapo informacija} (1,-8.15);
  \draw[dfdflow] (5.4,-7.05) -- node[srautas, pos=0.416, anchor=west] {O7 Garso\\instrukcijos} (5.4,-8.15);
  \draw[dfdflow] (-2.37,-3.2) -- node[srautas, pos=0.608, sloped, above] {O3.6.2 Valdymo\\pulto duomenys} (-5.85,-2.7);
  \draw[dfdflow] (-2.37,-4) -- node[srautas, pos=0.543, sloped, above] {O5 Pranešimas apie\\ryšio sutrikimą} (-5.85,-4.5);
  \draw[dfdflow] (-1.55,-0.85) -- node[srautasv, pos=0.172, anchor=east] {V33.3} (-1.55,-2.56);
  \draw[dfdflow] (-1.25,-2.56) -- node[srautasv, pos=0.704, anchor=west] {V34.2} (-1.25,-0.85);
  \draw[dfdflow] (-3.6,-6.15) -- node[srautasv, pos=0.146, anchor=west] {V48} (-2.24,-4.23);
\end{tikzpicture}
\caption{2 lygio duomenų srautų diagrama: 6 procesas „Vykdyti užsiėmimą“}\label{fig:dfd-2-6}
\end{figure}
% <<< GEN:dfd-2-6

\clearpage
\subsubsection{Duomenų srautų aprašai}\label{sec:dfd-srautai}

Srautų identifikatoriai: I~– įeinantys į sistemą, O~– išeinantys iš jos,
V~– vidiniai (tarp procesų ir saugyklų). Taškas rodo išskaidymą: I1.2 yra I1
dalis, o I1.2.1~– I1.2 dalis. Apibendrintų kontekstinės diagramos srautų
(\ref{tab:dfd-0}~lentelė) perduodami duomenys, pavidalas, intensyvumas ir
perdavimo sąlygos nurodyti juos sudarančių 1~lygio srautų aprašuose. Išoriniai srautai aprašyti
\ref{tab:dfd-io}~lentelėje, vidiniai 1~lygio srautai~– \ref{tab:dfd-v},
o 2~lygio srautai~– \ref{tab:dfd-l2}~lentelėje. Dauguma vidinių srautų yra
duomenų bazės įrašai, perduodami kartu su juos sukėlusiu išoriniu srautu;
išimtys nurodytos \ref{tab:dfd-v} ir \ref{tab:dfd-l2}~lentelėse. Stulpelyje
„FP funkcijos“ nurodytos projekto plano A~priedo funkcijos, todėl kiekvieną
srautą galima susieti su funkcinių taškų skaičiavimu.

Intensyvumas įvertintas vienam sporto centrui su baseinu ir sporto sale:
apie 10~užsiėmimų per dieną ir iki 20~dalyvių grupėje, todėl per dieną gali
būti iki 200~rezervacijų ir tiek pat priminimų. Šias prielaidas reikės
patikslinti su užsakovu.

\begingroup
\small
\begin{longtable}{@{}>{\raggedright\arraybackslash}p{0.9cm}>{\raggedright\arraybackslash}p{5.0cm}>{\raggedright\arraybackslash}p{2.4cm}>{\raggedright\arraybackslash}p{2.4cm}>{\raggedright\arraybackslash}p{3.6cm}@{}}
\caption{Kontekstinės diagramos srautų sudėtis}\label{tab:dfd-0}\\
\toprule
\textbf{ID} & \textbf{Pavadinimas} & \textbf{Šaltinis} & \textbf{Adresatas} & \textbf{1 lygio srautai} \\
\midrule
\endfirsthead
\toprule
\textbf{ID} & \textbf{Pavadinimas} & \textbf{Šaltinis} & \textbf{Adresatas} & \textbf{1 lygio srautai} \\
\midrule
\endhead
% >>> GEN:ctx
I1 & Dalyvio duomenys ir veiksmai & Dalyvis & Sistema & I1.1, I1.2 \\
O1 & Informacija dalyviui & Sistema & Dalyvis & O1.1, O1.2, O1.3 \\
O2 & Pranešimai dalyviui & Sistema & Dalyvis & O2.1, O2.2, O2.3 \\
I2 & Trenerio duomenys ir komandos & Treneris & Sistema & I2.1, I2.2, I2.3, I2.4, I2.5, I2.6 \\
O3 & Informacija treneriui & Sistema & Treneris & O3.1, O3.2, O3.3, O3.4, O3.5, O3.6, O3.7 \\
I3 & Administravimo duomenys ir turinys & Administratorius & Sistema & I3.1, I3.2, I3.3, I3.4 \\
O4 & Administravimo informacija & Sistema & Administratorius & O4.1, O4.2, O4.3, O4.4, O4.5 \\
O5 & Pranešimas apie ryšio sutrikimą & Sistema & Administratorius & tas pats srautas 1 lygyje \\
O6 & Pratimo vaizdas ir etapo informacija & Sistema & Ekranas & tas pats srautas 1 lygyje \\
O7 & Garso instrukcijos & Sistema & Garso įranga & tas pats srautas 1 lygyje \\
% <<< GEN:ctx
\bottomrule
\end{longtable}
\endgroup

\begin{landscape}
\begingroup
\footnotesize
\setlength{\tabcolsep}{3pt}
\begin{longtable}{@{}>{\raggedright\arraybackslash}p{0.9cm}>{\raggedright\arraybackslash}p{2.3cm}>{\raggedright\arraybackslash}p{4.4cm}>{\raggedright\arraybackslash}p{2.2cm}>{\raggedright\arraybackslash}p{2.2cm}>{\raggedright\arraybackslash}p{2.5cm}>{\raggedright\arraybackslash}p{3.0cm}>{\raggedright\arraybackslash}p{3.3cm}>{\raggedright\arraybackslash}p{1.7cm}@{}}
\caption{Išoriniai duomenų srautai}\label{tab:dfd-io}\\
\toprule
\textbf{ID} & \textbf{Pavadinimas} & \textbf{Perduodami duomenys} & \textbf{Šaltinis} & \textbf{Adresatas} & \textbf{Pavidalas} & \textbf{Intensyvumas} & \textbf{Perdavimo sąlygos} & \textbf{FP funkcijos} \\
\midrule
\endfirsthead
\caption[]{Išoriniai duomenų srautai (tęsinys)}\\
\toprule
\textbf{ID} & \textbf{Pavadinimas} & \textbf{Perduodami duomenys} & \textbf{Šaltinis} & \textbf{Adresatas} & \textbf{Pavidalas} & \textbf{Intensyvumas} & \textbf{Perdavimo sąlygos} & \textbf{FP funkcijos} \\
\midrule
\endhead
% >>> GEN:io
I1.1 & Dalyvio paskyros duomenys & Registracijos duomenys (vardas, pavardė, gimimo data, el. paštas, telefonas, plaukimo lygis, slaptažodis); prisijungimo el. paštas ir slaptažodis; slaptažodžio atkūrimo užklausa; pakeisti profilio duomenys, įskaitant tikslus ir pastabas treneriui & Dalyvis & 1 procesas & Žiniatinklio formos duomenys & Registracija~– kartą kiekvienam dalyviui; prisijungimas~– kaskart pradedant darbą & Registruojamas el. paštas dar nenaudojamas; profilį keičia prisijungęs dalyvis & P0.4–P0.6, P0.8 \\ \addlinespace[3pt]
O1.1 & Dalyvio profilis & Prisijungimo rezultatas; slaptažodžio atkūrimo nuoroda; dalyvio vardas, pavardė, gimimo data, el. paštas, telefonas, plaukimo lygis, tikslai ir pastabos treneriui & 1 procesas & Dalyvis & Žiniatinklio puslapis; atkūrimo nuoroda~– el. laišku & Kaskart prisijungus ar atidarius profilį & Rodomi tik paties dalyvio duomenys & P0.5, P0.7 \\ \addlinespace[3pt]
I1.2 & Rezervacijos veiksmas & Pasirinktas užsiėmimas ir veiksmas: rezervuoti vietą arba atšaukti savo rezervaciją & Dalyvis & 4 procesas & Žiniatinklio sąsajos veiksmas & Iki 200 rezervacijų ir atšaukimų per dieną & Dalyvis prisijungęs, jo paskyra aktyvi; užsiėmimas dar neprasidėjo ir neatšauktas & P3.3, P3.4 \\ \addlinespace[3pt]
O1.2 & Tvarkaraštis ir rezervacijos & Užsiėmimų tvarkaraštis (pavadinimas, režimas, data ir laikas, sporto erdvė, treneris, laisvų vietų skaičius); rezervacijos patvirtinimas arba atmetimo priežastis; dalyvio rezervacijos su būsenomis; lankomumo istorija & 4 procesas & Dalyvis & Žiniatinklio puslapis, pritaikytas telefonui & Kaskart atidarius tvarkaraštį ar savo rezervacijas ir po kiekvieno rezervacijos veiksmo & Pateikiama pagal dalyvio pasirinktą laikotarpį ir sporto erdvę; rodomi tik paties dalyvio duomenys & P3.2, P3.3, P3.8, P3.10 \\ \addlinespace[3pt]
O1.3 & Asmeninė programa & Aktuali asmeninė programa: pratimai, pasiruošimo rekomendacijos, namų užduotys & 5 procesas & Dalyvis & Žiniatinklio puslapis, pritaikytas telefonui & Kaskart dalyviui atidarius programą & Dalyviui paskirta asmeninė programa & P4.10 \\ \addlinespace[3pt]
O2.1 & Užsiėmimo pakeitimo pranešimas & Užsiėmimo pavadinimas ir pakeisti duomenys (data, laikas, sporto erdvė) arba žinia apie atšaukimą & 3 procesas & Dalyvis & El. laiškas & Kaskart pakeitus ar atšaukus užsiėmimą~– kiekvienam užsiregistravusiam (iki 20 laiškų) & Dalyvis turi aktyvią to užsiėmimo rezervaciją & P2.4, P2.5 \\ \addlinespace[3pt]
O2.2 & Priminimas apie užsiėmimą & Dalyvio vardas, užsiėmimo pavadinimas, data, laikas ir sporto erdvė & 4 procesas & Dalyvis & El. laiškas & Kartą per dieną nustatytu laiku~– kiekvienai rytdienos aktyviai rezervacijai (iki 200 laiškų) & Rezervacija aktyvi, užsiėmimas neatšauktas, priminimas dar nesiųstas & P3.11 \\ \addlinespace[3pt]
O2.3 & Programos pakeitimo pranešimas & Žinia, kad treneris pakeitė asmeninę programą ar namų užduotis & 5 procesas & Dalyvis & El. laiškas & Kaskart treneriui pakeitus programą & Pakeista dalyviui paskirta galiojanti programa & P4.4 \\ \addlinespace[3pt]
I2.1 & Trenerio paskyros duomenys & Prisijungimo el. paštas ir slaptažodis; slaptažodžio atkūrimo užklausa; pakeisti profilio duomenys & Treneris & 1 procesas & Žiniatinklio formos duomenys & Kaskart pradedant darbą & Trenerio paskyrą sukūrė administratorius & P0.5, P0.6, P0.8 \\ \addlinespace[3pt]
O3.1 & Trenerio profilis & Prisijungimo rezultatas; slaptažodžio atkūrimo nuoroda; trenerio profilio duomenys & 1 procesas & Treneris & Žiniatinklio puslapis; atkūrimo nuoroda~– el. laišku & Kaskart prisijungus ar atidarius profilį & Rodomi tik paties trenerio duomenys & P0.5, P0.7 \\ \addlinespace[3pt]
I2.2 & Programos duomenys & Programos pavadinimas ir aprašas; etapai (eilės numeris, trukmė, instrukcijos tekstas, pratimai iš katalogo, neprivaloma garso instrukcija iš mokomosios medžiagos); šalinimo ar archyvavimo nurodymas & Treneris & 2 procesas & Žiniatinklio formos duomenys & Kelis kartus per savaitę; pradiniame turinyje apie 10 programų & Treneris prisijungęs; kataloge yra įkeltų pratimų ir medžiagos & P1.10–P1.15 \\ \addlinespace[3pt]
O3.2 & Pratimai, medžiaga ir programos & Pratimų katalogas su demonstracija; mokomosios medžiagos sąrašas ir įrašų peržiūra; programų sąrašas, programos etapai ir versijų istorija & 2 procesas & Treneris & Žiniatinklio puslapis su vaizdo peržiūra ir garso perklausa & Kaskart sudarant ar peržiūrint programą & Treneris prisijungęs & P1.16–P1.22 \\ \addlinespace[3pt]
I2.3 & Užsiėmimo duomenys & Pavadinimas, data, pradžios laikas, trukmė, sporto erdvė, programa, tipas (individualus ar grupinis), vietų skaičius, vaizdo ir garso nustatymai (jei nenurodyti~– iš erdvės profilio); pakeitimai; atšaukimo nurodymas & Treneris & 3 procesas & Žiniatinklio formos duomenys & Iki 10 naujų užsiėmimų per dieną; pakeitimai pavieniai & Treneris tvarko tik savo užsiėmimus; jo užsiėmimų laikas nesikerta; keisti ir atšaukti galima tik dar neprasidėjusį užsiėmimą & P2.2, P2.4, P2.5 \\ \addlinespace[3pt]
O3.3 & Užsiėmimų sąrašas & Trenerio vedami ir virtualaus trenerio užsiėmimai su užimtumu (užsiregistravusių dalyvių ir laisvų vietų skaičius); pasirinkto užsiėmimo informacija & 3 procesas & Treneris & Žiniatinklio puslapis & Kaskart atidarius sąrašą & Rodomi trenerio vedami ir virtualaus trenerio užsiėmimai & P2.11, P2.13 \\ \addlinespace[3pt]
I2.4 & Priskyrimas ir lankomumas & Priskiriamo dalyvio el. paštas ir užsiėmimas arba šalinamas dalyvis; lankomumo žymos „dalyvavo“ ir „nedalyvavo“ & Treneris & 4 procesas & Žiniatinklio formos duomenys & Lankomumas~– kartą per užsiėmimą (iki 10 per dieną); priskyrimai~– keli per dieną & Treneris valdo tą užsiėmimą; priskirti ir pašalinti galima, kol užsiėmimas nėra atšauktas, užbaigtas ar neįvykęs; lankomumas žymimas prasidėjusiam ir neatšauktam užsiėmimui & P3.5–P3.7 \\ \addlinespace[3pt]
O3.4 & Dalyvių sąrašai & Priskyrimo ir lankomumo išsaugojimo rezultatai; užsiregistravusių dalyvių sąrašas; trenerio dalyvių sąrašas ir paieška; dalyvio kortelė; dalyvio lankomumo istorija & 4 procesas & Treneris & Žiniatinklio puslapis & Kaskart atidarius sąrašą ar kortelę & Rodomi tik trenerio vedamų ir virtualaus trenerio užsiėmimų dalyviai & P3.9, P3.10, P3.12, P3.13 \\ \addlinespace[3pt]
I2.5 & Įvertinimas ir asmeninė programa & Individualus dalyvio įvertinimas; įvertinimas, pagal kurį sudaroma programa; asmeninės programos pratimai ir pasiruošimo rekomendacijos; namų užduotys; pakeitimai, užbaigimo ar archyvavimo nurodymas & Treneris & 5 procesas & Žiniatinklio formos duomenys & Keli įvertinimai ir programų pakeitimai per savaitę vienam treneriui; programą galima keisti nuotoliniu būdu bet kada & Dalyvis turi ar turėjo rezervaciją užsiėmime, kurį veda šis treneris (virtualaus trenerio užsiėmimai neįskaičiuojami) & P4.2–P4.8 \\ \addlinespace[3pt]
O3.5 & Asmeninės programos & Trenerio sudarytų asmeninių programų sąrašas; pasirinkta programa su namų užduotimis & 5 procesas & Treneris & Žiniatinklio puslapis & Kaskart atidarius sąrašą ar programą & Rodomos tik trenerio sudarytos programos & P4.9, P4.10 \\ \addlinespace[3pt]
I2.6 & Vykdymo komandos & Pradėti, pristabdyti, tęsti, praleisti etapą, pereiti prie pasirinkto etapo, pakeisti etapo trukmę, užbaigti, keisti instrukcijų garsumą ar jas nutildyti & Treneris & 6 procesas & Valdymo pulto komanda & Nuo kelių iki kelių dešimčių komandų per užsiėmimą, įskaitant virtualaus trenerio užsiėmimus & Užsiėmimą valdo jį vedantis treneris, o virtualaus trenerio užsiėmimą~– bet kuris prisijungęs treneris; turinys parengtas grotuve & P2.6–P2.10, P6.2 \\ \addlinespace[3pt]
O3.6 & Valdymo pulto duomenys & Komandos rezultatas; dabartinis ir kitas etapas, likęs etapo laikas, programos eiga, grotuvo ryšio būsena & 6 procesas & Treneris & Valdymo pultas, atnaujinamas realiuoju laiku & Nuolat užsiėmimo metu, įskaitant virtualaus trenerio užsiėmimus & Užsiėmimas pradėtas & P2.12 \\ \addlinespace[3pt]
O3.7 & Istorija ir ataskaitos & Užbaigtų, neįvykusių ir atšauktų užsiėmimų, kuriuos treneris vedė ar valdė, istorija su faktine eiga ir lankomumo ataskaita (užsiregistravusių ir dalyvavusių skaičius, lankomumo procentas~– atšauktiems neskaičiuojamas, nepažymėtų skaičius) & 7 procesas & Treneris & Žiniatinklio puslapis & Kartą per savaitę ar mėnesį, pasirinkus laikotarpį & Rodomi tik užsiėmimai, kuriuos treneris vedė ar valdė & P7.1, P7.2 \\ \addlinespace[3pt]
I3.1 & Paskyros, erdvės ir įranga & Prisijungimo duomenys; slaptažodžio atkūrimo užklausa; pakeisti administratoriaus paskyros duomenys; naudotojo paskyros duomenys ir vaidmuo; deaktyvavimo arba aktyvavimo nurodymas; sporto erdvės duomenys ir jos vaizdo bei garso profilis; įrenginio registracija, nustatymai ir konfigūracijos pasas (programinės įrangos versija, IP adresas, būsena); šalinimo nurodymai & Administratorius & 1 procesas & Žiniatinklio formos duomenys & Prisijungimas~– kaskart pradedant darbą; kiti veiksmai~– kelis kartus per savaitę, daugiausia diegimo metu & Prisijungimo duomenys ir slaptažodžio atkūrimo užklausa priimami neprisijungus, kitiems veiksmams administratorius prisijungęs; šalinamoje erdvėje neturi būti jokių užsiėmimų (ir praėjusių ar atšauktų), o deaktyvuojamam treneriui~– būsimų užsiėmimų & P0.5, P0.6, P0.8, P0.9–P0.12, P0.15–P0.17, P0.19–P0.21 \\ \addlinespace[3pt]
O4.1 & Paskyrų, erdvių ir įrangos sąrašai & Prisijungimo rezultatas; slaptažodžio atkūrimo nuoroda; administratoriaus paskyros duomenys; naudotojų sąrašas su filtru pagal vaidmenį; paskyros informacija; erdvių sąrašas; erdvės įrangos sąrašas; įrenginio nustatymai & 1 procesas & Administratorius & Žiniatinklio puslapis; atkūrimo nuoroda~– el. laišku & Kaskart prisijungus ar atidarius sąrašą, įrašą ar savo paskyrą & Sąrašai ir paskyros duomenys~– prisijungusiam administratoriui; atkūrimo nuoroda siunčiama paskyros el. paštu & P0.5, P0.7, P0.13, P0.14, P0.18, P0.22, P0.23 \\ \addlinespace[3pt]
I3.2 & Pratimai ir medžiagos failai & Patvirtinti vaizdo ir garso failai su metaduomenimis; pratimo duomenys ir jo sąsaja su medžiaga; pakeitimai ir šalinimo nurodymai & Administratorius & 2 procesas & Žiniatinklio formos duomenys ir įkeliami failai & Pradinis įkėlimas~– apie 40 pratimų; vėliau pavieniai pakeitimai & Vaizdo failas neviršija Full HD raiškos; šalinamas pratimas nenaudojamas programose ir namų užduotyse, o įrašas~– pratimuose ir programose & P1.4–P1.9 \\ \addlinespace[3pt]
O4.2 & Pratimai, medžiaga ir programos & Pratimų katalogas su demonstracija; mokomosios medžiagos sąrašas, įrašų peržiūra ir perklausa; programų sąrašas & 2 procesas & Administratorius & Žiniatinklio puslapis su vaizdo peržiūra ir garso perklausa & Kaskart tvarkant turinį ar planuojant užsiėmimą & Administratorius prisijungęs & P1.16–P1.20 \\ \addlinespace[3pt]
I3.3 & Užsiėmimo duomenys & Užsiėmimas su pasirinktu treneriu arba virtualaus trenerio užsiėmimas: pavadinimas, programa, sporto erdvė, data, pradžios laikas, trukmė, tipas (individualus ar grupinis), vietų skaičius, vaizdo ir garso nustatymai (jei nenurodyti~– iš erdvės profilio); pakeitimai; atšaukimo nurodymas & Administratorius & 3 procesas & Žiniatinklio formos duomenys & Kartu su I2.3 iki 10 naujų užsiėmimų per dieną; pakeitimai ir atšaukimai pavieniai & Pasirinkto trenerio užsiėmimų laikas nesikerta; keisti ir atšaukti galima tik dar neprasidėjusį užsiėmimą & P2.2–P2.5 \\ \addlinespace[3pt]
O4.3 & Užsiėmimų sąrašas & Užsiėmimai su užimtumu (užsiregistravusių dalyvių ir laisvų vietų skaičius); pasirinkto užsiėmimo informacija & 3 procesas & Administratorius & Žiniatinklio puslapis & Kaskart atidarius sąrašą & Administratorius prisijungęs & P2.11, P2.13 \\ \addlinespace[3pt]
I3.4 & Dalyvių priskyrimas & Priskiriamo dalyvio el. paštas ir užsiėmimas arba šalinamas dalyvis & Administratorius & 4 procesas & Žiniatinklio formos duomenys & Keli priskyrimai per dieną & Užsiėmimas nėra atšauktas, užbaigtas ar neįvykęs; priskiriant~– yra laisvų vietų ir dalyvio paskyra aktyvi & P3.5, P3.6 \\ \addlinespace[3pt]
O4.4 & Užsiregistravę dalyviai & Priskyrimo rezultatas; užsiregistravusių dalyvių sąrašas (vardas, pavardė, telefonas, rezervacijos būsena, rezervavimo laikas) & 4 procesas & Administratorius & Žiniatinklio puslapis & Kaskart atidarius sąrašą & Administratorius prisijungęs & P3.9 \\ \addlinespace[3pt]
O4.5 & Istorija ir ataskaitos & Užbaigtų, neįvykusių ir atšauktų užsiėmimų istorija su faktine eiga; lankomumo ataskaita (įskaitant nepažymėtų skaičių); užsiėmimų ir erdvių užimtumo ataskaita, įskaitant virtualaus trenerio užsiėmimus & 7 procesas & Administratorius & Žiniatinklio puslapis & Kartą per savaitę ar mėnesį, pasirinkus laikotarpį & Administratorius prisijungęs & P7.1–P7.3 \\ \addlinespace[3pt]
O5 & Pranešimas apie ryšio sutrikimą & Užsiėmimo pavadinimas, sporto erdvė, sutrikimo laikas, vykdyto etapo numeris & 6 procesas & Administratorius & Pranešimas administratoriaus sąsajoje & Kaskart nutrūkus ryšiui su grotuvu; turėtų pasitaikyti retai & Užsiėmimas vykdomas arba pristabdytas, o serveris ilgiau už nustatytą ribą negauna grotuvo signalo & P2.14 \\ \addlinespace[3pt]
O6 & Pratimo vaizdas ir etapo informacija & Dabartinio etapo pratimų vaizdo įrašai eilės tvarka, instrukcijos tekstas, likęs laikas ir kitas etapas & 6 procesas & Ekranas & Full HD vaizdo signalas iš vietinio grotuvo & Nuolat užsiėmimo metu & Užsiėmimas vykdomas; turinys atsisiųstas į grotuvą & P5.1 \\ \addlinespace[3pt]
O7 & Garso instrukcijos & Įrašyta dabartinio etapo garso instrukcija nustatytu garsumu & 6 procesas & Garso įranga & Garso signalas iš vietinio grotuvo į radijo siųstuvą ar kolonėles & Nuolat užsiėmimo metu, sinchroniškai su vaizdu & Užsiėmimas vykdomas; instrukcijos nenutildytos & P6.1, P6.2 \\ \addlinespace[3pt]
% <<< GEN:io
\bottomrule
\end{longtable}

\begin{longtable}{@{}>{\raggedright\arraybackslash}p{0.9cm}>{\raggedright\arraybackslash}p{2.6cm}>{\raggedright\arraybackslash}p{2.3cm}>{\raggedright\arraybackslash}p{2.3cm}>{\raggedright\arraybackslash}p{5.8cm}>{\raggedright\arraybackslash}p{2.4cm}>{\raggedright\arraybackslash}p{3.4cm}>{\raggedright\arraybackslash}p{3.2cm}@{}}
\caption{Vidiniai duomenų srautai tarp procesų ir saugyklų}\label{tab:dfd-v}\\
\toprule
\textbf{ID} & \textbf{Pavadinimas} & \textbf{Šaltinis} & \textbf{Adresatas} & \textbf{Perduodami duomenys} & \textbf{Pavidalas} & \textbf{Intensyvumas} & \textbf{Perdavimo sąlygos} \\
\midrule
\endfirsthead
\caption[]{Vidiniai duomenų srautai tarp procesų ir saugyklų (tęsinys)}\\
\toprule
\textbf{ID} & \textbf{Pavadinimas} & \textbf{Šaltinis} & \textbf{Adresatas} & \textbf{Perduodami duomenys} & \textbf{Pavidalas} & \textbf{Intensyvumas} & \textbf{Perdavimo sąlygos} \\
\midrule
\endhead
% >>> GEN:v
V1 & Paskyros įrašas & 1 procesas & D1 Naudotojų paskyros & Nauja ar pakeista paskyra, vaidmuo ir būsena (aktyvi ar neaktyvi) & Duomenų bazės įrašas & Kartu su I1.1, I2.1, I3.1 & Kaip įeinančio srauto, procesui jį priėmus \\ \addlinespace[3pt]
V2 & Paskyrų duomenys & D1 Naudotojų paskyros & 1 procesas & Duomenys prisijungimui ir el. pašto unikalumui patikrinti, profilis, naudotojų sąrašas & Duomenų bazės įrašas & Kartu su I1.1, I2.1, I3.1, O1.1, O3.1, O4.1 & Kaip intensyvumo stulpelyje nurodytų srautų \\ \addlinespace[3pt]
V3 & Erdvės įrašas & 1 procesas & D2 Sporto erdvės & Nauja, pakeista ar pašalinta erdvė su vaizdo ir garso profiliu & Duomenų bazės įrašas & Kartu su I3.1 & Kaip įeinančio srauto, procesui jį priėmus \\ \addlinespace[3pt]
V4 & Erdvių duomenys & D2 Sporto erdvės & 1 procesas & Erdvių sąrašas ir erdvės duomenys & Duomenų bazės įrašas & Kartu su I3.1, O4.1 & Kaip intensyvumo stulpelyje nurodytų srautų \\ \addlinespace[3pt]
V5 & Įrenginio įrašas & 1 procesas & D3 Įranga & Nauja, pakeista ar pašalinta įrenginio registracija su nustatymais ir konfigūracijos pasu (programinės įrangos versija, IP adresas, būsena) & Duomenų bazės įrašas & Kartu su I3.1 & Kaip įeinančio srauto, procesui jį priėmus \\ \addlinespace[3pt]
V6 & Įrangos duomenys & D3 Įranga & 1 procesas & Erdvės įrangos sąrašas ir įrenginio nustatymai & Duomenų bazės įrašas & Kartu su I3.1, O4.1 & Kaip intensyvumo stulpelyje nurodytų srautų \\ \addlinespace[3pt]
V7 & Erdvės ir trenerio užsiėmimai & D7 Užsiėmimai & 1 procesas & Visi erdvės užsiėmimai (ir praėję ar atšaukti) bei būsimi trenerio užsiėmimai~– šalinimo ir deaktyvavimo patikrai & Duomenų bazės įrašas & Kaskart šalinant erdvę ar deaktyvuojant paskyrą & Administratorius šalina erdvę arba deaktyvuoja trenerio paskyrą \\ \addlinespace[3pt]
V8 & Rezervacijų atšaukimas & 1 procesas & D8 Rezervacijos ir lankomumas & Deaktyvuoto dalyvio būsimos rezervacijos pažymimos atšauktomis & Duomenų bazės įrašas & Kaskart deaktyvuojant dalyvio paskyrą & Deaktyvuojama dalyvio paskyra, turinti būsimų aktyvių rezervacijų \\ \addlinespace[3pt]
V9 & Pratimo įrašas & 2 procesas & D4 Pratimai & Naujas, pakeistas ar pašalintas pratimas ir jo sąsaja su medžiaga & Duomenų bazės įrašas & Kartu su I3.2 & Kaip įeinančio srauto, procesui jį priėmus \\ \addlinespace[3pt]
V10 & Pratimų duomenys & D4 Pratimai & 2 procesas & Pratimų katalogas, pratimo duomenys & Duomenų bazės įrašas & Kartu su I2.2, I3.2, O3.2, O4.2 & Kaip intensyvumo stulpelyje nurodytų srautų \\ \addlinespace[3pt]
V11 & Medžiagos įrašas & 2 procesas & D5 Mokomoji medžiaga & Įkeltas vaizdo ar garso failas su metaduomenimis, metaduomenų pakeitimai, pašalinamas įrašas & Vaizdo ar garso failas ir metaduomenys & Kartu su I3.2 & Kaip įeinančio srauto, procesui jį priėmus \\ \addlinespace[3pt]
V12 & Medžiagos duomenys & D5 Mokomoji medžiaga & 2 procesas & Medžiagos sąrašas, įrašai peržiūrai ir perklausai & Vaizdo ar garso failas ir metaduomenys & Kartu su I2.2, I3.2, O3.2, O4.2 & Kaip intensyvumo stulpelyje nurodytų srautų \\ \addlinespace[3pt]
V13 & Programos įrašas & 2 procesas & D6 Užsiėmimų programos & Nauja programa ar nauja jos versija su etapais, archyvavimo žyma & Duomenų bazės įrašas & Kartu su I2.2 & Kaip įeinančio srauto, procesui jį priėmus \\ \addlinespace[3pt]
V14 & Programų duomenys & D6 Užsiėmimų programos & 2 procesas & Programų sąrašas, etapai, versijų istorija, pratimų ir medžiagos naudojimas programose & Duomenų bazės įrašas & Kartu su I2.2, I3.2, O3.2, O4.2 & Kaip intensyvumo stulpelyje nurodytų srautų \\ \addlinespace[3pt]
V15 & Programų naudojimas & D7 Užsiėmimai & 2 procesas & Užsiėmimai, kuriems priskirta programa,~– šalinimo patikrai & Duomenų bazės įrašas & Kaskart šalinant ar archyvuojant programą & Treneris šalina ar archyvuoja programą \\ \addlinespace[3pt]
V16 & Pratimų naudojimas & D9 Asmeninės programos & 2 procesas & Asmeninės programos ir namų užduotys, kuriose naudojamas pratimas,~– šalinimo patikrai & Duomenų bazės įrašas & Kaskart šalinant pratimą & Administratorius šalina pratimą \\ \addlinespace[3pt]
V17 & Asmeninės programos įrašas & 5 procesas & D9 Asmeninės programos & Įvertinimas, nauja ar pakeista asmeninė programa su nuoroda į įvertinimą, pratimais ir pasiruošimo rekomendacijomis, namų užduotys, programos būsena (galiojanti, užbaigta ar archyvuota) & Duomenų bazės įrašas & Kartu su I2.5 & Kaip įeinančio srauto, procesui jį priėmus \\ \addlinespace[3pt]
V18 & Asmeninių programų duomenys & D9 Asmeninės programos & 5 procesas & Dalyvio programa, rekomendacijos, namų užduotys, trenerio programų sąrašas; dalyvio įvertinimai ir jais pagrįstos programos~– programos pagrindo patikrai & Duomenų bazės įrašas & Kartu su I2.5, O1.3, O3.5 & Kaip intensyvumo stulpelyje nurodytų srautų \\ \addlinespace[3pt]
V19 & Pratimai programai & D4 Pratimai & 5 procesas & Pratimai asmeninei programai ir namų užduotims sudaryti & Duomenų bazės įrašas & Kartu su I2.5, O1.3 & Kaip intensyvumo stulpelyje nurodytų srautų \\ \addlinespace[3pt]
V20 & Dalyvio duomenys & D1 Naudotojų paskyros & 5 procesas & Dalyvio vardas, pavardė ir el. paštas & Duomenų bazės įrašas & Kartu su I2.5, O2.3, O3.5 & Kaip intensyvumo stulpelyje nurodytų srautų \\ \addlinespace[3pt]
V21 & Užsiėmimo įrašas & 3 procesas & D7 Užsiėmimai & Naujas, pakeistas, atšauktas ar neįvykusiu pažymėtas užsiėmimas su programos versija ir turinio parengimo žyma & Duomenų bazės įrašas & Kartu su I2.3, I3.3; periodiškai~– neįvykusiems užsiėmimams pažymėti & Įrašoma, kai I2.3 ar I3.3 duomenys tenkina 3 proceso taisykles; periodiškai~– kai suplanuoto užsiėmimo pabaiga praėjo, o jis nepradėtas \\ \addlinespace[3pt]
V22 & Užsiėmimų duomenys & D7 Užsiėmimai & 3 procesas & Trenerio užsiėmimai laiko sankirtoms tikrinti, užsiėmimų sąrašas ir jų būsenos & Duomenų bazės įrašas & Kartu su I2.3, I3.3, O3.3, O4.3 & Kaip intensyvumo stulpelyje nurodytų srautų \\ \addlinespace[3pt]
V23 & Trenerių ir dalyvių duomenys & D1 Naudotojų paskyros & 3 procesas & Trenerių vardai; užsiregistravusių dalyvių vardai ir el. paštai pranešimams & Duomenų bazės įrašas & Kartu su I3.3, O2.1, O3.3, O4.3 & Kaip intensyvumo stulpelyje nurodytų srautų \\ \addlinespace[3pt]
V24 & Erdvių profiliai & D2 Sporto erdvės & 3 procesas & Sporto erdvės ir jų vaizdo bei garso profiliai užsiėmimo nustatymams & Duomenų bazės įrašas & Kartu su I2.3, I3.3, O3.3, O4.3 & Kaip intensyvumo stulpelyje nurodytų srautų \\ \addlinespace[3pt]
V25 & Programų duomenys & D6 Užsiėmimų programos & 3 procesas & Parengtos programos ir jų aktualios versijos užsiėmimui & Duomenų bazės įrašas & Kartu su I2.3, I3.3, O3.3, O4.3 & Kaip intensyvumo stulpelyje nurodytų srautų \\ \addlinespace[3pt]
V26 & Užsiregistravę dalyviai & D8 Rezervacijos ir lankomumas & 3 procesas & Aktyvios rezervacijos: užimtumas ir užsiregistravę dalyviai & Duomenų bazės įrašas & Kartu su O2.1, O3.3, O4.3 & Kaip intensyvumo stulpelyje nurodytų srautų \\ \addlinespace[3pt]
V27 & Rezervacijų atšaukimas & 3 procesas & D8 Rezervacijos ir lankomumas & Atšaukto užsiėmimo aktyvios rezervacijos pažymimos atšauktomis & Duomenų bazės įrašas & Kaskart atšaukus užsiėmimą & Atšaukiamas užsiėmimas, turintis aktyvių rezervacijų \\ \addlinespace[3pt]
V28 & Rezervacijos ir lankomumo įrašas & 4 procesas & D8 Rezervacijos ir lankomumas & Nauja, atšaukta ar pašalinta rezervacija, lankomumo žyma, žyma „priminimas išsiųstas“ & Duomenų bazės įrašas & Kartu su I1.2, I2.4, I3.4, O2.2 & Įrašoma, kai I1.2, I2.4 ar I3.4 duomenys tenkina 4 proceso taisykles; žyma „priminimas išsiųstas“~– sėkmingai išsiuntus O2.2 \\ \addlinespace[3pt]
V29 & Rezervacijų ir lankomumo duomenys & D8 Rezervacijos ir lankomumas & 4 procesas & Užimtos vietos, dalyvio rezervacijos, užsiregistravę dalyviai, lankomumo istorija & Duomenų bazės įrašas & Kartu su visais 4 proceso srautais & Kaip intensyvumo stulpelyje nurodytų srautų \\ \addlinespace[3pt]
V30 & Užsiėmimų duomenys & D7 Užsiėmimai & 4 procesas & Užsiėmimo pavadinimas, režimas, data ir laikas, vietų skaičius, būsena, treneris & Duomenų bazės įrašas & Kartu su visais 4 proceso srautais & Kaip intensyvumo stulpelyje nurodytų srautų \\ \addlinespace[3pt]
V31 & Dalyvių duomenys & D1 Naudotojų paskyros & 4 procesas & Paskyros būsena, vardas, pavardė, gimimo data, telefonas, el. paštas, plaukimo lygis, tikslai, pastabos treneriui; trenerio vardas & Duomenų bazės įrašas & Kartu su I1.2, I2.4, I3.4, O1.2, O2.2, O3.4, O4.4 & Kaip intensyvumo stulpelyje nurodytų srautų \\ \addlinespace[3pt]
V32 & Užsiėmimų sporto erdvės & D2 Sporto erdvės & 4 procesas & Sporto erdvė tvarkaraščiui ir priminimui & Duomenų bazės įrašas & Kartu su O1.2, O2.2 & Kaip intensyvumo stulpelyje nurodytų srautų \\ \addlinespace[3pt]
V33 & Vykdomas užsiėmimas & D7 Užsiėmimai & 6 procesas & Užsiėmimas (pavadinimas, vedantis treneris arba virtualaus trenerio režimas), jo programa ir priskirta versija, sporto erdvė, būsena, turinio parengimo ir ryšio žymos, vaizdo ir garso nustatymai & Duomenų bazės įrašas & Prieš užsiėmimą (6.1), kaskart gavus vykdymo komandą (6.2), grotuvo būseną ar nutrūkus ryšiui su grotuvu (6.6) & Užsiėmimas suplanuotas arba pradėtas \\ \addlinespace[3pt]
V34 & Faktinė eiga & 6 procesas & D7 Užsiėmimai & Užsiėmimo būsena, priskirta programos versija, turinio parengimo ir ryšio žymos, vykdymo žurnalo įrašai (faktinė etapų eiga, komandos ir jas davę treneriai, ryšio sutrikimai) & Duomenų bazės įrašas & Kaskart priėmus komandą, pasikeitus būsenai ar etapui, parengus turinį, nutrūkus ar atsikūrus ryšiui & Rengiant turinį grotuvui, priėmus leistiną trenerio komandą, grotuvui pranešus eigą arba nutrūkus ar atsikūrus ryšiui \\ \addlinespace[3pt]
V35 & Programos etapai & D6 Užsiėmimų programos & 6 procesas & Aktuali programos versija ir jos etapai: eiliškumas, trukmė, instrukcijos tekstas, pratimai, garso instrukcija & Duomenų bazės įrašas & Kartą prieš užsiėmimą ir pakartotinai iki jo pradžios & Užsiėmimas suplanuotas; turinys rengiamas grotuvui arba tikrinama, ar nesukurta naujesnė programos versija \\ \addlinespace[3pt]
V36 & Vaizdo ir garso įrašai & D5 Mokomoji medžiaga & 6 procesas & Programoje naudojami vaizdo ir garso failai & Vaizdo ir garso failai & Kartą prieš užsiėmimą ir iš naujo, jei turinys rengiamas pakartotinai & Prieš užsiėmimą, kai turinys atsiunčiamas į grotuvą \\ \addlinespace[3pt]
V37 & Įrenginių nustatymai & D3 Įranga & 6 procesas & Erdvės grotuvo, ekrano ir garso įrenginių nustatymai & Duomenų bazės įrašas & Kartą prieš užsiėmimą ir iš naujo, jei turinys rengiamas pakartotinai & Prieš užsiėmimą, kai turinys atsiunčiamas į grotuvą \\ \addlinespace[3pt]
V38 & Užsiėmimų istorija & D7 Užsiėmimai & 7 procesas & Užbaigti, neįvykę ir atšaukti užsiėmimai su vykdymo žurnalo įrašais (faktinė etapų eiga, komandas davę treneriai) & Duomenų bazės įrašas & Kartu su O3.7, O4.5 & Kaip intensyvumo stulpelyje nurodytų srautų \\ \addlinespace[3pt]
V39 & Rezervacijų ir lankomumo duomenys & D8 Rezervacijos ir lankomumas & 7 procesas & Rezervacijos ir lankomumo žymos ataskaitoms & Duomenų bazės įrašas & Kartu su O3.7, O4.5 & Kaip intensyvumo stulpelyje nurodytų srautų \\ \addlinespace[3pt]
V40 & Erdvių sąrašas & D2 Sporto erdvės & 7 procesas & Sporto erdvės užimtumo ataskaitai & Duomenų bazės įrašas & Kartu su O3.7, O4.5 & Kaip intensyvumo stulpelyje nurodytų srautų \\ \addlinespace[3pt]
V47 & Pratimai vykdymui & D4 Pratimai & 6 procesas & Programos etapuose naudojamų pratimų pavadinimai ir nuorodos į jų vaizdo įrašus & Duomenų bazės įrašas & Kartą prieš užsiėmimą ir iš naujo, jei turinys rengiamas pakartotinai & Prieš užsiėmimą, kai turinys atsiunčiamas į grotuvą \\ \addlinespace[3pt]
V48 & Sporto erdvė & D2 Sporto erdvės & 6 procesas & Užsiėmimo sporto erdvės pavadinimas pranešimui apie ryšio sutrikimą & Duomenų bazės įrašas & Kaskart nutrūkus ryšiui su grotuvu & Serveris ilgiau už nustatytą ribą negauna grotuvo signalo \\ \addlinespace[3pt]
V49 & Versijų autoriai & D1 Naudotojų paskyros & 2 procesas & Programos versijas sukūrusių trenerių vardai & Duomenų bazės įrašas & Kartu su O3.2 & Kaip intensyvumo stulpelyje nurodytų srautų \\ \addlinespace[3pt]
V50 & Trenerio užsiėmimai & D7 Užsiėmimai & 5 procesas & Užsiėmimai, kuriuos veda treneris (be virtualaus trenerio užsiėmimų),~– prieigos prie dalyvio patikrai & Duomenų bazės įrašas & Kartu su I2.5 & Kaip intensyvumo stulpelyje nurodytų srautų \\ \addlinespace[3pt]
V51 & Trenerio dalyvių rezervacijos & D8 Rezervacijos ir lankomumas & 5 procesas & Dalyvio rezervacijos trenerio užsiėmimuose~– prieigos prie dalyvio patikrai & Duomenų bazės įrašas & Kartu su I2.5 & Kaip intensyvumo stulpelyje nurodytų srautų \\ \addlinespace[3pt]
% <<< GEN:v
\bottomrule
\end{longtable}

\begin{longtable}{@{}>{\raggedright\arraybackslash}p{1.0cm}>{\raggedright\arraybackslash}p{1.15cm}>{\raggedright\arraybackslash}p{2.6cm}>{\raggedright\arraybackslash}p{2.6cm}>{\raggedright\arraybackslash}p{2.6cm}>{\raggedright\arraybackslash}p{6.6cm}>{\raggedright\arraybackslash}p{5.6cm}@{}}
\caption{2 lygio diagramų srautai}\label{tab:dfd-l2}\\
\toprule
\textbf{ID} & \textbf{Tėvinis} & \textbf{Pavadinimas} & \textbf{Šaltinis} & \textbf{Adresatas} & \textbf{Perduodami duomenys} & \textbf{Pavidalas, intensyvumas ir sąlygos} \\
\midrule
\endfirsthead
\caption[]{2 lygio diagramų srautai (tęsinys)}\\
\toprule
\textbf{ID} & \textbf{Tėvinis} & \textbf{Pavadinimas} & \textbf{Šaltinis} & \textbf{Adresatas} & \textbf{Perduodami duomenys} & \textbf{Pavidalas, intensyvumas ir sąlygos} \\
\midrule
\endhead
% >>> GEN:l2
I1.2.1 & I1.2 & Rezervavimo užklausa & Dalyvis & 4.2 procesas & Pasirinktas užsiėmimas ir veiksmas „Rezervuoti“ & Kaip I1.2 \\ \addlinespace[3pt]
I1.2.2 & I1.2 & Atšaukimo užklausa & Dalyvis & 4.3 procesas & Pasirinkta savo rezervacija ir veiksmas „Atšaukti“ & Kaip I1.2 \\ \addlinespace[3pt]
O1.2.1 & O1.2 & Tvarkaraštis & 4.1 procesas & Dalyvis & Užsiėmimai su pavadinimu, režimu, data ir laiku, sporto erdve, treneriu ir laisvų vietų skaičiumi & Kaip O1.2 \\ \addlinespace[3pt]
O1.2.2 & O1.2 & Rezervavimo rezultatas & 4.2 procesas & Dalyvis & Patvirtinimas arba atmetimo priežastis: nebėra vietų, užsiėmimas prasidėjo ar atšauktas, paskyra neaktyvi, rezervacija jau yra & Kaip O1.2 \\ \addlinespace[3pt]
O1.2.3 & O1.2 & Atšaukimo rezultatas & 4.3 procesas & Dalyvis & Patvirtinimas arba atmetimo priežastis: užsiėmimas jau prasidėjo & Kaip O1.2 \\ \addlinespace[3pt]
O1.2.4 & O1.2 & Rezervacijų istorija & 4.6 procesas & Dalyvis & Dalyvio rezervacijos su būsenomis ir lankomumo istorija & Kaip O1.2 \\ \addlinespace[3pt]
I2.4.1 & I2.4 & Dalyvio priskyrimas & Treneris & 4.4 procesas & Priskiriamo dalyvio el. paštas ir užsiėmimas arba šalinamas dalyvis & Kaip I2.4 \\ \addlinespace[3pt]
I2.4.2 & I2.4 & Lankomumo žymos & Treneris & 4.5 procesas & Užsiėmimo dalyvių žymos „dalyvavo“ ir „nedalyvavo“, jų pataisymai & Kaip I2.4 \\ \addlinespace[3pt]
O3.4.1 & O3.4 & Priskyrimo rezultatas & 4.4 procesas & Treneris & Patvirtinimas arba atmetimo priežastis: užsiėmimas atšauktas, užbaigtas ar neįvykęs, paskyra nerasta ar neaktyvi, nėra laisvų vietų, treneris užsiėmimo nevaldo, dalyvis jau užsiregistravęs & Kaip O3.4 \\ \addlinespace[3pt]
O3.4.2 & O3.4 & Lankomumo išsaugojimo rezultatas & 4.5 procesas & Treneris & Patvirtinimas arba atmetimo priežastis: užsiėmimas atšauktas ar dar neprasidėjo, treneris jo nevaldo & Kaip O3.4 \\ \addlinespace[3pt]
O3.4.3 & O3.4 & Dalyvių sąrašai & 4.6 procesas & Treneris & Užsiregistravusių dalyvių sąrašas, trenerio dalyvių sąrašas ir paieška, dalyvio kortelė, lankomumo istorija & Kaip O3.4 \\ \addlinespace[3pt]
O3.6.1 & O3.6 & Komandos rezultatas & 6.2 procesas & Treneris & Komanda priimta arba atmesta ir priežastis: turinys neparengtas, komanda neleistina esamoje būsenoje, nėra teisių & Kaip O3.6 \\ \addlinespace[3pt]
O3.6.2 & O3.6 & Valdymo pulto duomenys & 6.6 procesas & Treneris & Dabartinis ir kitas etapas, likęs etapo laikas, programos eiga, grotuvo ryšio būsena & Kaip O3.6 \\ \addlinespace[3pt]
O4.4.1 & O4.4 & Priskyrimo rezultatas & 4.4 procesas & Administratorius & Patvirtinimas arba atmetimo priežastis: užsiėmimas atšauktas, užbaigtas ar neįvykęs, paskyra nerasta ar neaktyvi, nėra laisvų vietų, dalyvis jau užsiregistravęs & Kaip O4.4 \\ \addlinespace[3pt]
O4.4.2 & O4.4 & Užsiregistravę dalyviai & 4.6 procesas & Administratorius & Užsiregistravusių dalyvių sąrašas (vardas, pavardė, telefonas, rezervacijos būsena, rezervavimo laikas) & Kaip O4.4 \\ \addlinespace[3pt]
V28.1 & V28 & Nauja rezervacija & 4.2 procesas & D8 Rezervacijos ir lankomumas & Aktyvi rezervacija: dalyvis, užsiėmimas, rezervavimo laikas, kas sukūrė rezervaciją & Kaip V28 \\ \addlinespace[3pt]
V28.2 & V28 & Atšaukta rezervacija & 4.3 procesas & D8 Rezervacijos ir lankomumas & Būsena „atšaukta“, atšaukimo laikas, kas atšaukė & Kaip V28 \\ \addlinespace[3pt]
V28.3 & V28 & Priskyrimas ar pašalinimas & 4.4 procesas & D8 Rezervacijos ir lankomumas & Nauja aktyvi rezervacija arba būsena „pašalinta“, veiksmo laikas ir jį atlikęs naudotojas & Kaip V28 \\ \addlinespace[3pt]
V28.4 & V28 & Lankomumo įrašas & 4.5 procesas & D8 Rezervacijos ir lankomumas & Lankomumo žyma, kas ją įrašė, įrašymo ir pataisymo laikas & Kaip V28 \\ \addlinespace[3pt]
V28.5 & V28 & Priminimo žyma & 4.7 procesas & D8 Rezervacijos ir lankomumas & Žyma „priminimas išsiųstas“ & Kaip V28 \\ \addlinespace[3pt]
V29.1 & V29 & Užimtos vietos & D8 Rezervacijos ir lankomumas & 4.1 procesas & Aktyvių rezervacijų skaičius kiekvienam užsiėmimui & Kaip V29 \\ \addlinespace[3pt]
V29.2 & V29 & Užimtos vietos & D8 Rezervacijos ir lankomumas & 4.2 procesas & Aktyvių rezervacijų skaičius ir dalyvio jau turima to užsiėmimo rezervacija & Kaip V29 \\ \addlinespace[3pt]
V29.3 & V29 & Dalyvio rezervacija & D8 Rezervacijos ir lankomumas & 4.3 procesas & Atšaukiamos rezervacijos savininkas ir būsena & Kaip V29 \\ \addlinespace[3pt]
V29.4 & V29 & Užimtos vietos & D8 Rezervacijos ir lankomumas & 4.4 procesas & Aktyvių rezervacijų skaičius ir dalyvio jau turima rezervacija & Kaip V29 \\ \addlinespace[3pt]
V29.5 & V29 & Užsiregistravę dalyviai & D8 Rezervacijos ir lankomumas & 4.5 procesas & Aktyvios užsiėmimo rezervacijos ir jau įrašytos lankomumo žymos & Kaip V29 \\ \addlinespace[3pt]
V29.6 & V29 & Rezervacijos ir lankomumas & D8 Rezervacijos ir lankomumas & 4.6 procesas & Rezervacijos su būsenomis ir rezervavimo laikais, lankomumo žymos & Kaip V29 \\ \addlinespace[3pt]
V29.7 & V29 & Rytdienos rezervacijos & D8 Rezervacijos ir lankomumas & 4.7 procesas & Rytdienos aktyvios rezervacijos be žymos „priminimas išsiųstas“ & Kaip V29 \\ \addlinespace[3pt]
V30.1 & V30 & Būsimi užsiėmimai & D7 Užsiėmimai & 4.1 procesas & Neatšaukti būsimi užsiėmimai: pavadinimas, režimas, data ir laikas, vietų skaičius & Kaip V30 \\ \addlinespace[3pt]
V30.2 & V30 & Užsiėmimo būsena & D7 Užsiėmimai & 4.2 procesas & Vietų skaičius; ar užsiėmimas dar neprasidėjo ir neatšauktas & Kaip V30 \\ \addlinespace[3pt]
V30.3 & V30 & Užsiėmimo pradžia & D7 Užsiėmimai & 4.3 procesas & Ar užsiėmimas dar neprasidėjo & Kaip V30 \\ \addlinespace[3pt]
V30.4 & V30 & Užsiėmimo duomenys & D7 Užsiėmimai & 4.4 procesas & Vietų skaičius, užsiėmimo būsena; ar užsiėmimą valdo šis treneris & Kaip V30 \\ \addlinespace[3pt]
V30.5 & V30 & Užsiėmimo duomenys & D7 Užsiėmimai & 4.5 procesas & Ar užsiėmimas prasidėjo ir neatšauktas; ar jį valdo šis treneris & Kaip V30 \\ \addlinespace[3pt]
V30.6 & V30 & Užsiėmimų duomenys & D7 Užsiėmimai & 4.6 procesas & Užsiėmimų pavadinimas, režimas, data ir laikas, treneris; treneriui~– tik jo vedami ir virtualaus trenerio užsiėmimai & Kaip V30 \\ \addlinespace[3pt]
V30.7 & V30 & Rytdienos užsiėmimai & D7 Užsiėmimai & 4.7 procesas & Pavadinimas, data ir laikas; ar užsiėmimas neatšauktas & Kaip V30 \\ \addlinespace[3pt]
V31.1 & V31 & Trenerio vardas & D1 Naudotojų paskyros & 4.1 procesas & Užsiėmimą vedančio trenerio vardas ir pavardė & Kaip V31 \\ \addlinespace[3pt]
V31.2 & V31 & Paskyros būsena & D1 Naudotojų paskyros & 4.2 procesas & Ar dalyvio paskyra aktyvi & Kaip V31 \\ \addlinespace[3pt]
V31.3 & V31 & Dalyvio paskyra & D1 Naudotojų paskyros & 4.4 procesas & Paskyra pagal el. paštą ir jos būsena & Kaip V31 \\ \addlinespace[3pt]
V31.4 & V31 & Dalyvių duomenys & D1 Naudotojų paskyros & 4.6 procesas & Vardas, pavardė, gimimo data, telefonas, el. paštas, plaukimo lygis, tikslai, pastabos treneriui; trenerio vardas & Kaip V31 \\ \addlinespace[3pt]
V31.5 & V31 & Gavėjo duomenys & D1 Naudotojų paskyros & 4.7 procesas & Dalyvio vardas ir el. paštas & Kaip V31 \\ \addlinespace[3pt]
V32.1 & V32 & Sporto erdvė & D2 Sporto erdvės & 4.1 procesas & Užsiėmimo sporto erdvė tvarkaraštyje & Kaip V32 \\ \addlinespace[3pt]
V32.2 & V32 & Sporto erdvė & D2 Sporto erdvės & 4.6 procesas & Užsiėmimo sporto erdvė rezervacijų sąraše & Kaip V32 \\ \addlinespace[3pt]
V32.3 & V32 & Sporto erdvė & D2 Sporto erdvės & 4.7 procesas & Užsiėmimo sporto erdvė priminime & Kaip V32 \\ \addlinespace[3pt]
V33.1 & V33 & Užsiėmimo nustatymai & D7 Užsiėmimai & 6.1 procesas & Užsiėmimas, jo programa ir priskirta versija, būsena, turinio parengimo žyma, erdvė ir vaizdo bei garso nustatymai & Kaip V33 \\ \addlinespace[3pt]
V33.2 & V33 & Užsiėmimo būsena & D7 Užsiėmimai & 6.2 procesas & Užsiėmimo būsena, turinio parengimo ir ryšio žymos, vedantis treneris arba virtualaus trenerio režimas & Kaip V33 \\ \addlinespace[3pt]
V33.3 & V33 & Vykdomo užsiėmimo duomenys & D7 Užsiėmimai & 6.6 procesas & Užsiėmimo pavadinimas, sporto erdvė ir būsena & Kaip V33 \\ \addlinespace[3pt]
V34.1 & V34 & Būsenos pakeitimas & 6.2 procesas & D7 Užsiėmimai & Nauja užsiėmimo būsena (vykdomas, pristabdytas, užbaigtas) ir vykdymo žurnalo įrašas apie gautą komandą ir ją davusį trenerį & Kaip V34 \\ \addlinespace[3pt]
V34.2 & V34 & Eiga ir ryšio būsena & 6.6 procesas & D7 Užsiėmimai & Būsenos pakeitimai pagal grotuvo signalą (be ryšio, užbaigtas pasibaigus paskutiniam etapui), žyma „ryšys atkurtas“ arba „ryšys neatkurtas“, vykdymo žurnalo įrašai apie faktinę etapų eigą ir ryšio sutrikimą & Kaip V34 \\ \addlinespace[3pt]
V34.3 & V34 & Parengimo žyma & 6.1 procesas & D7 Užsiėmimai & Priskirta programos versija ir žyma „turinys parengtas“ arba „turinys neparengtas“ & Kaip V34 \\ \addlinespace[3pt]
V41 & -- & Atsisiųstas turinys & 6.1 procesas & D10 Vietinė turinio kopija & Užsiėmimo programos versijos etapai, jų pratimai, pratimų demonstracijų vaizdo failai, etapų garso instrukcijų failai, užsiėmimo vaizdo ir garso nustatymai, įrenginių nustatymai & \textit{Pavidalas:} Failai grotuvo atmintyje. \textit{Intensyvumas:} Kartą prieš kiekvieną užsiėmimą ir iš naujo, jei iki pradžios pakeistas užsiėmimas ar jo programa. \textit{Sąlygos:} Užsiėmimas suplanuotas; grotuvas pasiekiamas tinklu. \\ \addlinespace[3pt]
V42 & -- & Etapo turinys & D10 Vietinė turinio kopija & 6.3 procesas & Dabartinio ir kito etapo vaizdas, garsas, instrukcijos tekstas ir trukmė & \textit{Pavidalas:} Failai grotuvo atmintyje. \textit{Intensyvumas:} Nuolat užsiėmimo metu. \textit{Sąlygos:} Turinys atsisiųstas; veikia ir nutrūkus ryšiui su serveriu. \\ \addlinespace[3pt]
V43 & -- & Vykdymo komanda & 6.2 procesas & 6.3 procesas & Patvirtinta komanda: pradėti, pristabdyti, tęsti, praleisti etapą, pereiti prie etapo, pakeisti trukmę, užbaigti, keisti garsumą ar nutildyti instrukcijas & \textit{Pavidalas:} Realiojo laiko pranešimas grotuvui. \textit{Intensyvumas:} Kaskart gavus leistiną komandą. \textit{Sąlygos:} Komandą davė užsiėmimą valdantis treneris. \\ \addlinespace[3pt]
V44 & -- & Etapo vaizdas & 6.3 procesas & 6.4 procesas & Dabartinio etapo pratimų vaizdo įrašai eilės tvarka, instrukcijos tekstas, likęs laikas, kitas etapas & \textit{Pavidalas:} Grotuvo vidiniai duomenys. \textit{Intensyvumas:} Nuolat užsiėmimo metu. \textit{Sąlygos:} Užsiėmimas vykdomas. \\ \addlinespace[3pt]
V45 & -- & Etapo garsas & 6.3 procesas & 6.5 procesas & Dabartinio etapo garso instrukcija, garsumo lygis ir nutildymo žyma & \textit{Pavidalas:} Grotuvo vidiniai duomenys. \textit{Intensyvumas:} Nuolat užsiėmimo metu, sinchroniškai su vaizdu. \textit{Sąlygos:} Užsiėmimas vykdomas. \\ \addlinespace[3pt]
V46 & -- & Grotuvo būsena & 6.3 procesas & 6.6 procesas & Dabartinio ir kito etapo eilės numeriai, likęs laikas, programos eiga ir grotuvo veikimo signalas & \textit{Pavidalas:} Periodinis būsenos pranešimas serveriui. \textit{Intensyvumas:} Kas kelias sekundes užsiėmimo metu; intervalą nustato nefunkciniai reikalavimai. \textit{Sąlygos:} Grotuvas pasiekiamas tinklu; nutrūkus ryšiui signalas negaunamas. \\ \addlinespace[3pt]
% <<< GEN:l2
\bottomrule
\end{longtable}
\endgroup
\end{landscape}

\subsubsection{Procesų specifikacijos}\label{sec:dfd-pspec}

Aprašyti visi žemiausio lygio procesai: neišskaidyti 1~lygio procesai 1, 2,
3, 5 ir 7 bei 4 ir 6~procesų dalys. Logika užrašyta struktūrizuota kalba su
raktiniais žodžiais \textbf{JEI}, \textbf{TADA}, \textbf{KITAIP},
\textbf{KIEKVIENAM} ir \textbf{KOL}. Prie kiekvieno proceso nurodyti jo
gaunami, skaitomi, rašomi ir pateikiami srautai. Rezervavimo (4.2) sąlygos,
nuo kurių priklauso rezultatas, papildomai pateiktos sprendimų lentelėje
(\ref{tab:pspec-42}~lentelė). Čia ir toliau „užsiėmimas prasidėjo“ reiškia,
kad atėjo suplanuotas jo pradžios laikas arba treneris jį jau pradėjo (6.2).

% >>> GEN:pspec
\paragraph{1 Tvarkyti paskyras, erdves ir įrangą}\mbox{}\par
{\small \textit{Gauna:} I1.1, I2.1, I3.1. \textit{Skaito:} V2, V4, V6, V7. \textit{Rašo:} V1, V3, V5, V8. \textit{Pateikia:} O1.1, O3.1, O4.1.}
\begin{itemize}
  \item Registracija: \textbf{JEI} trūksta privalomų laukų arba el. paštas jau yra D1 saugykloje, \textbf{TADA} atmesti ir nurodyti priežastį; \textbf{KITAIP} sukurti aktyvią dalyvio paskyrą.
  \item Prisijungimas: \textbf{JEI} paskyra yra, aktyvi ir slaptažodis teisingas, \textbf{TADA} suteikti prieigą pagal vaidmenį; \textbf{KITAIP} prieigos nesuteikti.
  \item Slaptažodžio atkūrimas: \textbf{JEI} el. paštas yra D1 saugykloje, \textbf{TADA} išsiųsti atkūrimo nuorodą; naują slaptažodį įrašyti į D1.
  \item Profilio, paskyros ir vaidmens pakeitimai: patikrinti laukus ir įrašyti į D1.
  \item Paskyros deaktyvavimas: \textbf{JEI} tai treneris ir D7 saugykloje yra jo būsimų užsiėmimų, \textbf{TADA} atmesti; \textbf{KITAIP} paskyrą pažymėti neaktyvia, istorijos netrinti, o \textbf{JEI} tai dalyvis, \textbf{TADA} jo būsimas aktyvias rezervacijas D8 saugykloje pažymėti atšauktomis. Aktyvavimas: neaktyvią paskyrą vėl pažymėti aktyvia.
  \item Erdvės šalinimas: \textbf{JEI} D7 saugykloje yra šios erdvės užsiėmimų (ir praėjusių ar atšauktų) arba D3 saugykloje yra jai priskirtos įrangos, \textbf{TADA} atmesti; \textbf{KITAIP} ištrinti.
  \item Įrenginio registravimas, nustatymų ir konfigūracijos paso keitimas, šalinimas: patikrinti nustatymus ir įrašyti į D3.
  \item Profilis ir sąrašai pateikiami tik pagal naudotojo vaidmenį.
\end{itemize}
\paragraph{2 Tvarkyti pratimus, medžiagą ir programas}\mbox{}\par
{\small \textit{Gauna:} I2.2, I3.2. \textit{Skaito:} V10, V12, V14, V15, V16, V49. \textit{Rašo:} V9, V11, V13. \textit{Pateikia:} O3.2, O4.2.}
\begin{itemize}
  \item Failo įkėlimas: \textbf{JEI} vaizdo raiška viršija Full HD arba failas netinkamo formato, \textbf{TADA} jo nepriimti ir nurodyti priežastį; \textbf{KITAIP} įrašyti failą ir metaduomenis į D5.
  \item Pratimo kūrimas ir keitimas: įrašyti pratimą į D4 kartu su nuoroda į medžiagą D5.
  \item Pratimo šalinimas: \textbf{JEI} pratimas naudojamas D6 saugyklos programose arba D9 saugyklos asmeninėse programose ir namų užduotyse, \textbf{TADA} atmesti; \textbf{KITAIP} ištrinti.
  \item Įrašo šalinimas: \textbf{JEI} įrašas naudojamas pratimuose ar programose, \textbf{TADA} atmesti; \textbf{KITAIP} ištrinti.
  \item Programos kūrimas: įrašyti į D6 kaip pirmąją versiją. Keitimas: sukurti naują versiją, nurodant ją sukūrusį trenerį; ankstesnė versija lieka istorijoje. Versijų istorijoje rodomi trenerių vardai iš D1.
  \item Etapų pridėjimas, keitimas ir šalinimas: patikrinti, ar trukmė teigiama ir ar eilės numeriai nesikartoja.
  \item Programos šalinimas: \textbf{JEI} programa priskirta D7 saugyklos užsiėmimams, \textbf{TADA} ją archyvuoti; \textbf{KITAIP} ištrinti. Treneris programą gali archyvuoti ir pats; archyvuota programa nebeparenkama naujiems užsiėmimams.
\end{itemize}
\paragraph{3 Planuoti užsiėmimus}\mbox{}\par
{\small \textit{Gauna:} I2.3, I3.3. \textit{Skaito:} V22, V23, V24, V25, V26. \textit{Rašo:} V21, V27. \textit{Pateikia:} O2.1, O3.3, O4.3.}
\begin{itemize}
  \item Planavimas: patikrinti privalomus laukus. \textbf{JEI} užsiėmimas individualus, \textbf{TADA} vietų skaičių nustatyti lygų 1. Nenurodytus vaizdo ir garso nustatymus perimti iš sporto erdvės profilio D2. Užsiėmimą susieti su aktualia pasirinktos programos versija iš D6; vykdomą versiją priskiria 6.1~procesas, rengdamas turinį grotuvui.
  \item \textbf{JEI} užsiėmimą veda treneris ir jo laikas kertasi su kitu to trenerio užsiėmimu D7 saugykloje, \textbf{TADA} atmesti ir nurodyti, su kuriuo užsiėmimu laikas sutampa.
  \item \textbf{JEI} tai virtualaus trenerio užsiėmimas, \textbf{TADA} trenerio nepriskirti ir pažymėti virtualaus trenerio režimą.
  \item Užsiėmimą įrašyti į D7 su būsena „suplanuotas“ ir žyma „turinys neparengtas“.
  \item Keitimas ir atšaukimas: \textbf{JEI} užsiėmimo būsena ne „suplanuotas“ arba jau atėjo jo pradžios laikas, \textbf{TADA} atmesti. Keičiant taikyti tas pačias taisykles; \textbf{JEI} naujas vietų skaičius mažesnis už aktyvių rezervacijų skaičių D8 saugykloje, \textbf{TADA} atmesti; \textbf{KITAIP} įrašyti pakeitimą ir žymą „turinys neparengtas“, kad 6.1~procesas turinį parengtų iš naujo.
  \item Periodiškai \textbf{KIEKVIENAM} užsiėmimui, kurio būsena „suplanuotas“, o suplanuota pabaiga jau praėjo: pakeisti būseną į „neįvykęs“.
  \item Atšaukimas: užsiėmimo būseną pakeisti į „atšauktas“, visas jo aktyvias rezervacijas D8 saugykloje pažymėti atšauktomis.
  \item Pakeitus ar atšaukus: \textbf{KIEKVIENAM} dalyviui, turinčiam aktyvią rezervaciją, išsiųsti O2.1 jo el. pašto adresu iš D1.
  \item Sąraše užimtumas = aktyvių rezervacijų skaičius, laisvos vietos = vietų skaičius $-$ užimtumas.
\end{itemize}
\paragraph{4.1 Pateikti tvarkaraštį}\mbox{}\par
{\small \textit{Skaito:} V29.1, V30.1, V31.1, V32.1. \textit{Pateikia:} O1.2.1.}
\begin{itemize}
  \item \textbf{KIEKVIENAM} pasirinkto laikotarpio ir sporto erdvės neatšauktam būsimam užsiėmimui: laisvos vietos = vietų skaičius $-$ aktyvių rezervacijų skaičius D8.
  \item Pateikti pavadinimą, režimą, datą ir laiką, sporto erdvę, trenerį ir laisvų vietų skaičių.
\end{itemize}
\paragraph{4.2 Rezervuoti vietą}\mbox{}\par
{\small \textit{Gauna:} I1.2.1. \textit{Skaito:} V29.2, V30.2, V31.2. \textit{Rašo:} V28.1. \textit{Pateikia:} O1.2.2.}
\begin{itemize}
  \item Patikrinti sąlygas pagal sprendimų lentelę.
  \item Laisvų vietų patikra ir rezervacijos įrašymas atliekami kaip vienas nedalomas veiksmas: kai paskutinę vietą vienu metu rezervuoja keli dalyviai, ją gauna tik vienas, kitiems pateikiamas pranešimas „Laisvų vietų nebėra“.
  \item \textbf{JEI} išsaugoti nepavyko, \textbf{TADA} patvirtinimo nerodyti.
\end{itemize}
\begin{table}[H]
\centering
\caption{Proceso 4.2 „Rezervuoti vietą“ sprendimų lentelė}\label{tab:pspec-42}
\small
\begin{tabular}{@{}lccccc@{}}
\toprule
\textbf{Sąlygos} & \textbf{R1} & \textbf{R2} & \textbf{R3} & \textbf{R4} & \textbf{R5} \\
\midrule
S1. Dalyvio paskyra aktyvi & T & N & T & T & T \\
S2. Užsiėmimas neprasidėjo ir neatšauktas & T & -- & N & T & T \\
S3. Dalyvis neturi aktyvios šio užsiėmimo rezervacijos & T & -- & -- & N & T \\
S4. Yra laisvų vietų & T & -- & -- & -- & N \\
\midrule
\textbf{Veiksmai} & & & & & \\
Sukurti rezervaciją ir parodyti patvirtinimą & X & & & & \\
Atmesti: „Paskyra neaktyvi“ & & X & & & \\
Atmesti: „Užsiėmimas prasidėjo arba atšauktas“ & & & X & & \\
Parodyti esamą rezervaciją, naujos nekurti & & & & X & \\
Atmesti: „Laisvų vietų nebėra“ & & & & & X \\
\bottomrule
\end{tabular}

\smallskip
{\footnotesize T~– taip, N~– ne, „--“~– nesvarbu.}
\end{table}
\paragraph{4.3 Atšaukti rezervaciją}\mbox{}\par
{\small \textit{Gauna:} I1.2.2. \textit{Skaito:} V29.3, V30.3. \textit{Rašo:} V28.2. \textit{Pateikia:} O1.2.3.}
\begin{itemize}
  \item \textbf{JEI} rezervacija priklauso dalyviui, yra aktyvi ir užsiėmimas dar neprasidėjo, \textbf{TADA} būseną pakeisti į „atšaukta“, įrašyti atšaukimo laiką ir naudotoją, pateikti patvirtinimą.
  \item \textbf{JEI} užsiėmimas jau prasidėjo, \textbf{TADA} atmesti; \textbf{JEI} rezervacija jau atšaukta ar pašalinta, \textbf{TADA} parodyti esamą būseną.
\end{itemize}
\paragraph{4.4 Priskirti ar pašalinti dalyvį}\mbox{}\par
{\small \textit{Gauna:} I2.4.1, I3.4. \textit{Skaito:} V29.4, V30.4, V31.3. \textit{Rašo:} V28.3. \textit{Pateikia:} O3.4.1, O4.4.1.}
\begin{itemize}
  \item \textbf{JEI} užsiėmimas atšauktas, užbaigtas arba neįvykęs, \textbf{TADA} priskyrimą ir pašalinimą atmesti.
  \item Priskyrimas: rasti paskyrą pagal el. paštą D1 saugykloje. \textbf{JEI} paskyra nerasta arba neaktyvi, treneris to užsiėmimo nevaldo, nėra laisvų vietų arba dalyvis jau turi aktyvią rezervaciją, \textbf{TADA} atmesti; \textbf{KITAIP} sukurti aktyvią rezervaciją ir įrašyti, kas ją sukūrė.
  \item Kaip ir 4.2 procese, laisvų vietų patikra ir rezervacijos įrašymas atliekami kaip vienas nedalomas veiksmas: jei paskutinę vietą vienu metu bando užimti keli dalyviai, ją gauna tik vienas.
  \item Pašalinimas: \textbf{JEI} rezervacija aktyvi ir užsiėmimą valdo šis treneris (administratoriui~– bet kurį užsiėmimą), \textbf{TADA} rezervacijos būseną pakeisti į „pašalinta“, įrašyti laiką ir naudotoją, vieta atlaisvinama; \textbf{KITAIP} atmesti.
  \item Rezultatą pateikti treneriui (O3.4.1) arba administratoriui (O4.4.1).
\end{itemize}
\paragraph{4.5 Žymėti lankomumą}\mbox{}\par
{\small \textit{Gauna:} I2.4.2. \textit{Skaito:} V29.5, V30.5. \textit{Rašo:} V28.4. \textit{Pateikia:} O3.4.2.}
\begin{itemize}
  \item \textbf{JEI} užsiėmimas prasidėjo, neatšauktas ir jį valdo šis treneris, \textbf{TADA} \textbf{KIEKVIENAM} užsiregistravusiam dalyviui įrašyti žymą „dalyvavo“ arba „nedalyvavo“, ją įrašiusį trenerį ir laiką; \textbf{KITAIP} atmesti. Lankomumą galima žymėti ir tada, kai užsiėmimas sistemoje nebuvo pradėtas.
  \item Taisant žymą įrašyti naują reikšmę ir pataisymo laiką. Kol žyma neįrašyta, lankomumas laikomas nepažymėtu, o ne „nedalyvavo“.
  \item Rezultatą pateikti treneriui (O3.4.2).
\end{itemize}
\paragraph{4.6 Pateikti rezervacijas ir dalyvių sąrašus}\mbox{}\par
{\small \textit{Skaito:} V29.6, V30.6, V31.4, V32.2. \textit{Pateikia:} O1.2.4, O3.4.3, O4.4.2.}
\begin{itemize}
  \item Dalyviui pateikti tik jo rezervacijas su būsenomis ir jo pasirinkto laikotarpio lankomumo istoriją.
  \item Treneriui pateikti tik jo vedamų ir virtualaus trenerio užsiėmimų dalyvius, dalyvio kortelę ir lankomumo istoriją; administratoriui~– bet kurio užsiėmimo užsiregistravusius dalyvius.
\end{itemize}
\paragraph{4.7 Siųsti priminimus}\mbox{}\par
{\small \textit{Skaito:} V29.7, V30.7, V31.5, V32.3. \textit{Rašo:} V28.5. \textit{Pateikia:} O2.2.}
\begin{itemize}
  \item Kasdien nustatytu laiku \textbf{KIEKVIENAM} rytdienos aktyviai rezervacijai be žymos „priminimas išsiųstas“, kurios užsiėmimas neatšauktas: parengti laišką (vardas, data, laikas, užsiėmimas, sporto erdvė) ir jį išsiųsti.
  \item \textbf{JEI} laiškas išsiųstas, \textbf{TADA} įrašyti žymą „priminimas išsiųstas“; \textbf{KITAIP} žymos neįrašyti. Pakartotinio siuntimo tvarką nustato nefunkciniai reikalavimai.
\end{itemize}
\paragraph{5 Tvarkyti asmenines programas}\mbox{}\par
{\small \textit{Gauna:} I2.5. \textit{Skaito:} V18, V19, V20, V50, V51. \textit{Rašo:} V17. \textit{Pateikia:} O1.3, O2.3, O3.5.}
\begin{itemize}
  \item \textbf{JEI} dalyvis neturi ir neturėjo rezervacijos užsiėmime, kurį veda šis treneris (tikrinama D8 ir D7 saugyklose; virtualaus trenerio užsiėmimai neįskaičiuojami), \textbf{TADA} atmesti; \textbf{KITAIP} įrašyti dalyvio įvertinimą į D9.
  \item \textbf{JEI} pasirinktas įvertinimas jau yra kitos programos pagrindas (D9), \textbf{TADA} atmesti; \textbf{KITAIP} sudaryti programą pagal jį: parinkti pratimus iš D4, pridėti pasiruošimo rekomendacijas ir įrašyti į D9 kaip galiojančią dalyvio programą.
  \item Keičiant programą ar namų užduotis įrašyti pakeitimą; \textbf{JEI} pakeista galiojanti programa, \textbf{TADA} išsiųsti dalyviui O2.3.
  \item Užbaigta ar archyvuota programa dalyviui nebepateikiama kaip aktuali, bet lieka istorijoje.
  \item Dalyviui pateikiama tik jo aktuali programa, treneriui~– tik jo sudarytos programos.
\end{itemize}
\paragraph{6.1 Parengti turinį grotuvui}\mbox{}\par
{\small \textit{Skaito:} V33.1, V35, V36, V37, V47. \textit{Rašo:} V34.3, V41.}
\begin{itemize}
  \item Prieš užsiėmimą priskirti jam tuo metu aktualią programos versiją iš D6 (ji pakeičia versiją, susietą planuojant), nuskaityti užsiėmimo nustatymus, tos versijos etapus, jų pratimus su demonstracijomis, etapų garso instrukcijas bei sporto erdvės įrangos nustatymus ir įrašyti juos į vietinę kopiją D10. Pradėjus užsiėmimą (6.2) versija nebekeičiama.
  \item \textbf{JEI} atsisiųsti pavyko, \textbf{TADA} D7 saugykloje įrašyti priskirtą versiją ir žymą „turinys parengtas“; \textbf{KITAIP}~– „turinys neparengtas“.
  \item \textbf{JEI} iki užsiėmimo pradžios jis pakeistas (žyma „turinys neparengtas“) arba D6 saugykloje sukurta naujesnė jo programos versija, \textbf{TADA} turinį parengti iš naujo.
\end{itemize}
\paragraph{6.2 Priimti vykdymo komandas}\mbox{}\par
{\small \textit{Gauna:} I2.6. \textit{Skaito:} V33.2. \textit{Rašo:} V34.1. \textit{Pateikia:} O3.6.1, V43.}
\begin{itemize}
  \item \textbf{JEI} komandą davė užsiėmimą valdantis treneris ir komanda leistina esamoje užsiėmimo būsenoje (žr. \ref{sec:gyvavimo-ciklai}~poskyrį), \textbf{TADA} perduoti ją grotuvui ir įrašyti į vykdymo žurnalą kartu su ją davusiu treneriu; \textbf{KITAIP} atmesti.
  \item Leidžiamos komandos: vykdomame užsiėmime~– visos; pristabdytame~– tęsti, užbaigti ir keisti garsumą; vykdomame be ryšio~– tęsti (tik atkūrus ryšį) ir užbaigti. „Pradėti“ leidžiama tik suplanuotam užsiėmimui, kai turinys parengtas.
  \item Užsiėmimo būseną D7 saugykloje keičia tik komandos „pradėti“, „pristabdyti“, „tęsti“ ir „užbaigti“; kitos komandos tik perduodamos grotuvui.
  \item Komandos rezultatą pateikti treneriui (O3.6.1).
\end{itemize}
\paragraph{6.3 Vykdyti programos eigą}\mbox{}\par
{\small \textit{Gauna:} V43. \textit{Skaito:} V42. \textit{Pateikia:} V44, V45, V46.}
\begin{itemize}
  \item Gavus komandą „pradėti“, pradėti pirmąjį etapą iš D10.
  \item \textbf{KOL} užsiėmimas vykdomas: perduoti dabartinio etapo vaizdą 6.4~procesui ir garsą 6.5~procesui, skaičiuoti likusį laiką; pasibaigus etapo laikui pereiti prie kito.
  \item Komandos: pristabdyti ir tęsti~– sustabdyti ir vėl paleisti laikmatį; praleisti ar pereiti~– pradėti nurodytą etapą; keisti trukmę~– perskaičiuoti likusį laiką; keisti garsumą ar nutildyti instrukcijas~– perduoti 6.5~procesui.
  \item Periodiškai siųsti grotuvo būseną; pasibaigus paskutiniam etapui pranešti, kad programa baigta.
  \item \textbf{JEI} ryšys su serveriu nutrūko, \textbf{TADA} užbaigti dabartinį etapą iš D10 ir naujo nepradėti.
\end{itemize}
\paragraph{6.4 Rodyti etapą ekrane}\mbox{}\par
{\small \textit{Gauna:} V44. \textit{Pateikia:} O6.}
\begin{itemize}
  \item Ekrane eilės tvarka rodyti dabartinio etapo pratimų vaizdo įrašus, instrukcijos tekstą, likusį laiką ir kitą etapą; pereinant prie kito etapo vaizdą atnaujinti.
\end{itemize}
\paragraph{6.5 Perduoti garso instrukcijas}\mbox{}\par
{\small \textit{Gauna:} V45. \textit{Pateikia:} O7.}
\begin{itemize}
  \item Perduoti dabartinio etapo garso instrukciją sinchroniškai su vaizdu, taikant nustatytą garsumą; \textbf{JEI} instrukcijos nutildytos, \textbf{TADA} jos neperduoti.
\end{itemize}
\paragraph{6.6 Stebėti eigą ir ryšį su grotuvu}\mbox{}\par
{\small \textit{Gauna:} V46. \textit{Skaito:} V33.3, V48. \textit{Rašo:} V34.2. \textit{Pateikia:} O3.6.2, O5.}
\begin{itemize}
  \item Gavus grotuvo būseną atnaujinti trenerio valdymo pultą ir į vykdymo žurnalą D7 saugykloje įrašyti faktinę etapų eigą (etapo pradžia ir pabaiga).
  \item \textbf{JEI} grotuvas praneša, kad programa baigta, \textbf{TADA} užsiėmimo būseną pakeisti į „užbaigtas“.
  \item \textbf{JEI} vykdomo ar pristabdyto užsiėmimo grotuvo signalas negaunamas ilgiau už nustatytą ribą, \textbf{TADA} pažymėti, kad užsiėmimas vykdomas be ryšio ir ryšys neatkurtas, išsiųsti O5 budinčiam darbuotojui (užsiėmimas iš D7, sporto erdvė iš D2) ir įrašyti sutrikimą į vykdymo žurnalą. Ribą nustato nefunkciniai reikalavimai.
  \item \textbf{JEI} signalas vėl gaunamas, \textbf{TADA} D7 saugykloje pažymėti, kad ryšys atkurtas.
\end{itemize}
\paragraph{7 Pateikti istoriją ir ataskaitas}\mbox{}\par
{\small \textit{Skaito:} V38, V39, V40. \textit{Pateikia:} O3.7, O4.5.}
\begin{itemize}
  \item Pagal pasirinktą laikotarpį atrinkti užbaigtus, neįvykusius ir atšauktus užsiėmimus; atšauktiems lankomumo procentas neskaičiuojamas.
  \item Treneriui pateikiama tik užsiėmimų, kuriuos jis vedė ar valdė, istorija ir lankomumo ataskaita; administratoriui~– visų užsiėmimų istorija, lankomumo ir užimtumo ataskaitos.
  \item Istorija: užsiėmimai su faktine eiga, įskaitant virtualaus trenerio užsiėmimus.
  \item Lankomumo ataskaita: užsiregistravę = rezervacijos, aktyvios užsiėmimo pradžioje, ir rezervacijos, sukurtos jau prasidėjus užsiėmimui (dalyviui atėjus be rezervacijos); dalyvavę = žymos „dalyvavo“; lankomumo procentas = dalyvavę / užsiregistravę $\times$ 100. Nepažymėtas lankomumas rodomas atskirai.
  \item Užimtumo ataskaita: kiekvienai sporto erdvei~– užsiėmimų skaičius (atšaukti nurodomi atskirai) ir vietų užimtumas = jos užsiėmimų aktyvios rezervacijos / tų užsiėmimų vietų skaičių suma $\times$ 100, atšauktų užsiėmimų neįskaičiuojant.
\end{itemize}
% <<< GEN:pspec
